\chapter{Flattening Stratification of a Coherent Sheaf}
\label{chap:Picard_FlatteningStratification}
% archon:covers AlgebraicJacobian/Picard/FlatteningStratification.lean


\section{Setup and motivation}
\label{sec:flatstrat_setup}

Generic flatness is the foundational input to the flattening-stratification
construction underlying Grothendieck's existence theorem for the Quot scheme
(\cite{nitsure-hilbert-quot}, \S4--\S5).  Its content is that flatness of a
coherent sheaf is a generic property over an integral base: over a noetherian
integral scheme \(S\), a coherent \(\OO_X\)-module on a finite-type morphism
\(X \to S\) becomes flat after restricting to a non-empty open subscheme of
\(S\). Algebraically this is the statement that a finite module over a
finite-type algebra over a noetherian domain becomes free after inverting a
single non-zero element of the base ring.

We work throughout with a noetherian (integral, where required) base scheme
\(S\), a finite-type morphism \(p : X \to S\), and a coherent \(\OO_X\)-module
\(\F\).

\section{Generic flatness}
\label{sec:generic_flatness}

The construction of the flattening stratification proceeds by Noetherian
induction on \(S\). The induction step relies on the following classical
\emph{generic flatness} theorem of Grothendieck: over a noetherian integral
base, a coherent sheaf is flat over a non-empty open subscheme. Algebraically,
this is the assertion that a finitely generated module over a finite-type
algebra over a noetherian domain is generically free after inverting a single
non-zero element of the base ring.

\begin{theorem}



\leanok
[Generic flatness, algebraic form]
  \label{thm:generic_flatness_algebraic}
  \lean{AlgebraicGeometry.genericFlatnessAlgebraic}
  \uses{lem:gf_torsion_base, lem:noeth_prime_filtration, lem:gf_splice_shortExact,
        lem:gf_noether_clear_denominators, lem:gf_polynomial_core, lem:gf_finite_module}
  % SOURCE: [Nitsure], \S4 ``Lemma on Generic Flatness''
  % (read from references/nitsure-hilbert-quot-src/nitsure-hilbert-quot.tex, L1711-L1716).
  % SOURCE QUOTE:
  % "Let \(A\) be a noetherian domain,
  %  and \(B\) a finite type \(A\) algebra. Let \(M\) be a finite \(B\)-module.
  %  Then there exists an \(f\in A\), \(f\ne 0\), such that
  %  the localisation \(M_f\) is a free module over \(A_f\)."
  \textit{Source: [Nitsure], \S4, ``Lemma on Generic Flatness''.}
  Let \(A\) be a noetherian domain, \(B\) a finite-type \(A\)-algebra, and \(M\) a finite
  \(B\)-module. Then there exists \(f \in A\), \(f \ne 0\), such that the localisation
  \(M_f\) is a free \(A_f\)-module.

\end{theorem}

% SOURCE QUOTE PROOF: "Over the quotient field \(K\) of \(A\), the \(K\)-algebra
% \(B_K = K\otimes _AB\) is of finite type, and \(M_K = K\otimes_AM\) is a
% finite module over \(B_K\). Let \(n\) be the dimension of the support
% of \(M_K\) over \(\spec B_K\). We argue by induction on \(n\), starting
% with \(n=-1\) which is the case when \(M_K =0\). [...]
% Now let \(n\ge 0\), and let the lemma be proved for smaller values.
% As \(B\) is noetherian and \(M\) is assumed to be a finite \(B\)-module,
% there exists a finite filtration
% \(0=M_0 \subset \ldots  \subset M_r =M\)
% where each \(M_i\) is a \(B\)-submodule of \(M\) such that for each \(i\ge 1\)
% the quotient module \(M_i/M_{i-1}\) is isomorphic to \(B/\pp_i\) for
% some prime ideal \(\pp_i\) in \(B\). [...] By Noether normalisation lemma,
% there exist elements \(b_1,\ldots,b_n \in B\), such that \(K\otimes_AB\) is
% finite over its subalgebra \(K[b_1,\ldots, b_n]\) [...]"
% (Source: nitsure-hilbert-quot-src/nitsure-hilbert-quot.tex, L1719-L1772; full
% proof omitted for length, structural steps restated below in project notation.)
\begin{proof}


\leanok
  \textbf{Primary route — the finite \(A\)-module case.}
  Suppose first that \(M\) is module-finite over \(A\) itself. Since \(A\) is noetherian,
  \(M\) is finitely presented over \(A\). The generic fibre \(M_K = K \otimes_A M\)
  (where \(K = \mathrm{Frac}(A)\)) is a finite-dimensional \(K\)-vector space, hence free
  over \(K\); by \cref{lem:gf_finite_module} this freeness descends to a single basic open:
  there exists \(f \in A\), \(f \ne 0\), with \(M_f\) free over \(A_f\). This closes the
  theorem whenever \(M\) is finite over \(A\).

  \textbf{Surviving residue — reducing the finite-type \(B\)-algebra case to the
  finite \(A\)-module case.} The genuine content not supplied by Mathlib is the
  reduction of the general hypothesis (\(M\) finite over a \emph{finite-type}
  \(A\)-algebra \(B\)) to the finite \(A\)-module situation handled above. Concretely,
  the missing piece is the polynomial-ring core of generic freeness: a finite module
  over \(A[b_1, \dots, b_n]\) becomes free over \(A_f\) after inverting a single
  \(f \in A\), \(f \ne 0\). The classical \S4 dévissage below is the fallback route
  for exactly this residue.

  \emph{Fallback (Nitsure \S4 induction).}
  The \S4 dévissage is assembled from the sub-lemmas of
  \cref{sec:gf_devissage_chain}. If \(M_K = 0\), the torsion base case
  \cref{lem:gf_torsion_base} gives \(f \ne 0\) with \(M_f = 0\) free. Otherwise the
  prime-filtration induction principle \cref{lem:noeth_prime_filtration} over the
  noetherian ring \(B\), combined with the splicing fact
  \cref{lem:gf_splice_shortExact}, reduces to \(M = B/\mathfrak{p}\) and hence to
  \(B\) a domain, \(M = B\); Noether normalisation with clearing of denominators
  \cref{lem:gf_noether_clear_denominators} makes \(B_g\) finite over a polynomial
  ring \(A_g[b_1, \dots, b_n]\), and the polynomial-ring core
  \cref{lem:gf_polynomial_core} — bottoming out at the finite-\(A\)-module leaf
  \cref{lem:gf_finite_module} of the primary route — produces the required
  \(f \ne 0\) with \(M_f\) free over \(A_f\).
\end{proof}

\subsection{Dévissage chain for the finite-type residue}
\label{sec:gf_devissage_chain}

The following sub-lemmas decompose the surviving residue of
\cref{thm:generic_flatness_algebraic} — the case where \(M\) is finite over the
finite-type \(A\)-algebra \(B\) but not module-finite over \(A\) — into the
classical \cite{nitsure-hilbert-quot}~\S4 induction. They are ordered so that each block
depends only on earlier blocks or Mathlib lemmas. Throughout,
\(A\) is a noetherian domain with fraction field \(K = \mathrm{Frac}(A)\), \(B\) a
finite-type \(A\)-algebra, and \(M\) a finite \(B\)-module with the compatible
\(A\)-module structure via the scalar tower \(A \to B \to M\).

\begin{lemma}\leanok
[Finite-module leaf — finite \(A\)-module case]
  \label{lem:gf_finite_module}
  \lean{AlgebraicGeometry.GenericFreeness.exists_free_localizationAway_of_finite}
  \uses{lem:fp_free_descent}
  Let \(A\) be a noetherian domain and \(M\) a finite \(A\)-module. Then there
  exists \(f \in A\), \(f \ne 0\), such that \(M_f\) is free over \(A_f\).
\end{lemma}

\begin{proof}


\leanok
  This is the \(d = 0\) leaf of the induction. A finite module over a noetherian ring is
  finitely presented, so
  \(M\) is finitely presented over \(A\). Over the generic point
  \(K = \mathrm{Frac}(A) = S^{-1}A\) (\(S = A \setminus \{0\}\)) the localisation
  \(M_K = S^{-1}M\) is a finite-dimensional \(K\)-vector space, hence free over
  \(K\). The finite-presentation free-descent lemma \cref{lem:fp_free_descent}
  then descends this freeness to a single basic open: there is \(r \in S\), i.e.\
  \(r \in A\) with \(r \ne 0\), such that \(M_r\) is free over \(A_r\). Take
  \(f := r\).
\end{proof}

\begin{lemma}[Finite-presentation free descent]
  \label{lem:fp_free_descent}
  \lean{Module.FinitePresentation.exists_free_localizedModule_powers}
  \mathlibok
  \textit{Provided by Mathlib.}
  Let \(R\) be a commutative ring, \(S \subseteq R\) a multiplicative subset, and
  \(M\) a finitely-presented \(R\)-module such that the localisation \(M_S\) is
  free over the localisation \(R_S\). Then there exists \(r \in S\) such that
  \(M_r\) is free over \(R_r\) (of the same finite rank).
\end{lemma}

\begin{lemma}\leanok
[L1 — torsion base case \(M_K = 0\)]
  \label{lem:gf_torsion_base}
  \lean{AlgebraicGeometry.GenericFreeness.exists_free_localizationAway_of_torsion}
  % SOURCE: [Nitsure], \S4 ``Lemma on Generic Flatness'', proof, base case
  % (read from references/nitsure-hilbert-quot-src/nitsure-hilbert-quot.tex, L1719-L1726).
  % SOURCE QUOTE PROOF:
  % "We argue by induction on $n$, starting with $n=-1$ which is the case when
  %  $M_K =0$. In this case, as $K\otimes _AM = S^{-1}M$ where $S = A - \{ 0\}$,
  %  each $v\in M$ is annihilated by some non-zero element of $A$. Taking a finite
  %  generating set, and a common multiple of corresponding annihilating elements,
  %  we see there exists an $f\ne 0$ in $A$ with $fM=0$. Hence $M_f =0$, proving
  %  the lemma when $n=-1$."
  \textit{Source: [Nitsure], \S4, ``Lemma on Generic Flatness'' (base case of the
  induction).}
  Suppose \(K \otimes_A M = 0\), where \(K = \mathrm{Frac}(A)\). Then there exists
  \(f \in A\), \(f \ne 0\), with \(f M = 0\); consequently \(M_f = 0\) is free over
  \(A_f\).
\end{lemma}

\begin{proof}


\leanok
  % SOURCE QUOTE PROOF: see the comment above the statement (L1719-L1726).
  \textit{Source: [Nitsure], \S4 (base case).}
  Since \(K \otimes_A M = S^{-1}M\) with \(S = A \setminus \{0\}\), the hypothesis
  \(M_K = 0\) means every element of \(M\) is killed by some non-zero element of
  \(A\). Choose a finite \(B\)-generating set \(v_1, \dots, v_k\) of \(M\) (which
  exists because \(M\) is a finite \(B\)-module), and for each \(i\) an
  \(a_i \in A\), \(a_i \ne 0\), with \(a_i v_i = 0\). The product
  \(f := a_1 \cdots a_k\) is non-zero (\(A\) is a domain) and annihilates every
  generator, hence \(f M = 0\). Therefore \(M_f = 0\), which is free over \(A_f\).
\end{proof}

\begin{lemma}[L2 — prime-filtration induction principle]
  \label{lem:noeth_prime_filtration}
  \lean{IsNoetherianRing.induction_on_isQuotientEquivQuotientPrime}
  \mathlibok
  \textit{Provided by Mathlib.}
  Let \(B\) be a noetherian commutative ring and let \(P\) be a predicate on finite
  \(B\)-modules that (i) holds for every subsingleton module, (ii) holds for every
  cyclic quotient \(B/\mathfrak{p}\) with \(\mathfrak{p}\) prime, and (iii) is
  closed under short exact sequences (\(P(N_1)\) and \(P(N_3)\) imply \(P(N_2)\)
  for \(0 \to N_1 \to N_2 \to N_3 \to 0\)). Then \(P\) holds for every finite
  \(B\)-module. In particular every finite \(B\)-module \(M\) admits a finite
  filtration \(0 = M_0 \subset \cdots \subset M_r = M\) with each quotient
  \(M_i / M_{i-1} \cong B/\mathfrak{p}_i\) for a prime \(\mathfrak{p}_i\), so to
  prove generic freeness for \(M\) it suffices to prove it for each
  \(B/\mathfrak{p}\) and to splice along the filtration's short exact sequences.
\end{lemma}

\begin{lemma}\leanok
[L3a — localisation of a short exact sequence is exact]
  \label{lem:gf_splice_shortExact_localized_exact}
  \lean{AlgebraicGeometry.GenericFreeness.exact_localizedModule_powers_of_shortExact}
  Let \(0 \to M' \xrightarrow{i} M \xrightarrow{q} M'' \to 0\) be a short exact
  sequence of \(B\)-modules, and let \(f \in A\), \(f \ne 0\). Restricting scalars
  along \(A \to B\) and localising every term at the powers of \(f\) yields a short
  exact sequence
  \[
    0 \to M'_f \xrightarrow{i_f} M_f \xrightarrow{q_f} M''_f \to 0
  \]
  of \(A_f\)-modules; in particular \(i_f\) is injective and \(q_f\) is surjective.
\end{lemma}

\begin{proof}


\leanok
  Localisation at a multiplicative set is an exact functor, preserving injections,
  surjections, and exactness. Applying it to the scalar-restricted \(A\)-linear maps
  \(i\) and \(q\) (via \(A \to B\)) returns the displayed sequence.
\end{proof}

\begin{lemma}\leanok
[L3b — free transport across a finer localisation]
  \label{lem:gf_splice_shortExact_free_transport}
  \lean{AlgebraicGeometry.GenericFreeness.free_localizationAway_of_free_of_eq_mul}
  Let \(N\) be a \(B\)-module and \(f', f'' \in A\) with \(f = f' f''\). If the
  localisation \(N_{f'}\) is free over \(A_{f'}\), then the further localisation
  \(N_f\) is free over \(A_f\).
\end{lemma}

\begin{proof}


\leanok
  Because \(f = f' f''\), the localisation away from \(f\) factors through the
  localisation away from \(f'\): the ring \(A_f\) is the localisation of \(A_{f'}\)
  at the image of \(f''\), and likewise the module \(N_f\) is the localisation of
  \(N_{f'}\) at the image of \(f''\). A localisation of a free module is free over
  the localised base ring — a basis of \(N_{f'}\) over \(A_{f'}\) maps to a basis of
  \(N_f\) over \(A_f\). Hence \(N_f\) is free over \(A_f\). Stated as a standalone
  fact, it is applied to both ends of the spliced sequence.
\end{proof}

\begin{lemma}\leanok
[L3c — a short exact sequence with free ends has a free middle]
  \label{lem:gf_splice_shortExact_split}
  \lean{AlgebraicGeometry.GenericFreeness.free_of_shortExact_of_free_free}
  Let \(0 \to M'_f \to M_f \to M''_f \to 0\) be a short exact sequence of
  \(A_f\)-modules in which both ends \(M'_f\) and \(M''_f\) are free over \(A_f\).
  Then \(M_f\) is free over \(A_f\); concretely \(M_f \cong M'_f \oplus M''_f\).
\end{lemma}

\begin{proof}


\leanok
  A free module is projective, so the free quotient \(M''_f\) is projective and the
  short exact sequence splits: \(M_f \cong M'_f \oplus M''_f\). A finite direct sum
  of free \(A_f\)-modules is free, hence \(M_f\) is free over \(A_f\).
\end{proof}

\begin{lemma}\leanok
[L3 — splicing fact for a short exact sequence (assembly)]
  \label{lem:gf_splice_shortExact}
  \lean{AlgebraicGeometry.GenericFreeness.exists_free_localizationAway_of_shortExact}
  \uses{lem:gf_splice_shortExact_localized_exact,
        lem:gf_splice_shortExact_free_transport, lem:gf_splice_shortExact_split}
  % SOURCE: [Nitsure], \S4 ``Lemma on Generic Flatness'', proof, splicing remark
  % (read from references/nitsure-hilbert-quot-src/nitsure-hilbert-quot.tex, L1737-L1745).
  % SOURCE QUOTE PROOF:
  % "Note that if $0\to M'\to M\to M'' \to 0$ is a short exact sequence of
  %  $B$-modules, and if $f'$ and $f''$ are non-zero elements of $A$ such that
  %  $M'_{f'}$ and $M''_{f''}$ are free over respectively $A_{f'}$ and $A_{f''}$,
  %  then $M_f$ is a free module over $A_f$ where $f=f'f''$. We will use this fact
  %  repeatedly. Therefore it is enough to prove the result when $M$ is of the form
  %  $B/\pp$ for a prime ideal $\pp$ in $B$. This reduces us to the case where $B$
  %  is a domain and $M=B$."
  \textit{Source: [Nitsure], \S4, ``Lemma on Generic Flatness'' (splicing fact).}
  Let \(0 \to M' \to M \to M'' \to 0\) be a short exact sequence of \(B\)-modules,
  and let \(f', f'' \in A\) be non-zero with \(M'_{f'}\) free over \(A_{f'}\) and
  \(M''_{f''}\) free over \(A_{f''}\). Then for \(f := f' f''\) the localisation
  \(M_f\) is free over \(A_f\).
\end{lemma}

\begin{proof}
  \uses{lem:gf_splice_shortExact_localized_exact,
        lem:gf_splice_shortExact_free_transport, lem:gf_splice_shortExact_split}
        \leanok
  % SOURCE QUOTE PROOF: see the comment above the statement (L1737-L1745).
  \textit{Source: [Nitsure], \S4 (splicing fact).}
  Set \(f := f' f''\); since \(A\) is a domain and \(f', f'' \ne 0\), the product is
  non-zero, so \(f \ne 0\). Localising the given short exact sequence at the powers
  of \(f\) (after restricting scalars to \(A\)) produces a short exact sequence
  \(0 \to M'_f \to M_f \to M''_f \to 0\) of \(A_f\)-modules, by
  \cref{lem:gf_splice_shortExact_localized_exact}. Both ends are free over \(A_f\):
  applying \cref{lem:gf_splice_shortExact_free_transport} to \(M'\) with the
  factorisation \(f = f' f''\) (using that \(M'_{f'}\) is free over \(A_{f'}\)) shows
  \(M'_f\) is free over \(A_f\), and applying the same lemma to \(M''\) with the
  factorisation \(f = f'' f'\) (using that \(M''_{f''}\) is free over \(A_{f''}\))
  shows \(M''_f\) is free over \(A_f\). Finally
  \cref{lem:gf_splice_shortExact_split} concludes that the middle term \(M_f\) is
  free over \(A_f\).
\end{proof}

\begin{lemma}[L4 anchor — Noether normalisation over a field]
  \label{lem:noether_normalization_fg}
  \lean{exists_finite_inj_algHom_of_fg}
  \mathlibok
  \textit{Provided by Mathlib.}
  Let \(k\) be a field and \(R\) a non-trivial finite-type \(k\)-algebra. Then there
  exist \(s \in \mathbb{N}\) and an injective \(k\)-algebra homomorphism
  \(k[X_1, \dots, X_s] \to R\) that is module-finite, i.e.\ \(R\) is a finite module
  over a polynomial subalgebra on algebraically independent generators.
\end{lemma}

\begin{lemma}


\leanok
[L4a — clearing one equation's denominators]
  \label{lem:gf_clear_one_denominator}
  \lean{AlgebraicGeometry.GenericFreeness.gf_clear_one_denominator}
  % SOURCE: [Nitsure], \S4 ``Lemma on Generic Flatness'', proof, normalisation step
  % (read from references/nitsure-hilbert-quot-src/nitsure-hilbert-quot.tex, L1755-L1759).
  % SOURCE QUOTE PROOF:
  % "If $g\ne 0$ in $A$ is chosen to be a `common denominator' for
  %  coefficients of equations of integral dependence
  %  satisfied by a finite set of
  %  algebra generators for $K\otimes_AB$ over $K[b_1,\ldots, b_n]$,
  %  we see that $B_g$ is finite over $A_g[b_1,\ldots, b_n]$."
  \textit{Source: [Nitsure], \S4, ``Lemma on Generic Flatness'' (common denominator).}
  Let \(A\) be a noetherian domain with fraction field \(K = \mathrm{Frac}(A)\), and let
  \(p \in K[X_1, \dots, X_n]\) be a polynomial with coefficients in \(K\) (for instance,
  one coefficient polynomial of an equation of integral dependence over
  \(K[b_1, \dots, b_n]\)). Then there exists \(g \in A\), \(g \ne 0\), such that \(p\)
  lies in the image of the coefficient-extension map
  \(A_g[X_1, \dots, X_n] \to K[X_1, \dots, X_n]\); equivalently, \(p\) has all of its
  coefficients in \(A_g = A[g^{-1}] \subseteq K\).
\end{lemma}

\begin{proof}

\leanok
  % SOURCE QUOTE PROOF: see the comment above the statement (L1755-L1759).
  \textit{Source: [Nitsure], \S4 (common denominator).}
  The polynomial \(p\) has finitely many non-zero coefficients
  \(c_1, \dots, c_r \in K = \mathrm{Frac}(A)\). Each \(c_j\) is a fraction
  \(c_j = a_j / s_j\) with \(a_j \in A\) and \(s_j \in A \setminus \{0\}\). Put
  \(g := s_1 \cdots s_r \in A\); since \(A\) is a domain and each \(s_j \ne 0\), also
  \(g \ne 0\). For each \(j\) the denominator \(s_j\) divides \(g\), so
  \(c_j = a_j / s_j\) lies in the subring \(A_g = A[g^{-1}] \subseteq K\). Hence every
  coefficient of \(p\) lies in \(A_g\), and \(p\) is the coefficient-extension of a
  polynomial in \(A_g[X_1, \dots, X_n]\).
\end{proof}

\begin{lemma}\leanok
[L4 — Noether normalisation with clearing of denominators]
  \label{lem:gf_noether_clear_denominators}
  \lean{AlgebraicGeometry.GenericFreeness.exists_localizationAway_finite_mvPolynomial}
  \uses{lem:noether_normalization_fg, lem:gf_clear_one_denominator,
        lem:gf_isLocalization_lift_injective}
  % SOURCE: [Nitsure], \S4 ``Lemma on Generic Flatness'', proof, normalisation step
  % (read from references/nitsure-hilbert-quot-src/nitsure-hilbert-quot.tex, L1747-L1758).
  % SOURCE QUOTE PROOF:
  % "By Noether normalisation lemma, there exist elements $b_1,\ldots,b_n \in B$,
  %  such that $K\otimes_AB$ is finite over its subalgebra $K[b_1,\ldots, b_n]$ and
  %  the elements $b_1,\ldots, b_n$ are algebraically independent over $K$. [...]
  %  If $g\ne 0$ in $A$ is chosen to be a `common denominator' for coefficients of
  %  equations of integral dependence satisfied by a finite set of algebra
  %  generators for $K\otimes_AB$ over $K[b_1,\ldots, b_n]$, we see that $B_g$ is
  %  finite over $A_g[b_1,\ldots, b_n]$."
  \textit{Source: [Nitsure], \S4, ``Lemma on Generic Flatness'' (Noether
  normalisation with denominators).}
  Let \(A\) be a noetherian domain with fraction field \(K\) and \(B\) a finite-type
  \(A\)-algebra that is a domain with \(B_K = K \otimes_A B \ne 0\). Write
  \(A_g = A[g^{-1}]\) and \(B_g = B[g^{-1}]\) for the localisations away from \(g\)
  and its image. Then there exist a natural number \(n\), an element \(g \in A\) with
  \(g \ne 0\), and an injective \(A_g\)-algebra homomorphism
  \[
    \varphi : A_g[X_1, \dots, X_n] \longrightarrow B_g,
  \]
  where \(A_g[X_1, \dots, X_n]\) is the polynomial ring on \(n\) algebraically
  independent variables, such that \(B_g\) is module-finite over the image polynomial
  subalgebra via \(\varphi\). The injectivity of \(\varphi\) records the algebraic
  independence of the generators \(b_j := \varphi(X_j)\), and module-finiteness is
  taken over the algebra structure that \(\varphi\) puts on \(B_g\).
\end{lemma}

\begin{proof}
  \uses{lem:noether_normalization_fg, lem:gf_clear_one_denominator}
  \leanok
  % SOURCE QUOTE PROOF: see the comment above the statement (L1747-L1758).
  \textit{Source: [Nitsure], \S4 (Noether normalisation with denominators).}
  \emph{Step 1 — Noether normalisation over \(K\).} Apply Noether normalisation over the
  field \(K\) (\cref{lem:noether_normalization_fg}) to the finite-type \(K\)-algebra
  \(B_K = K \otimes_A B\): there are algebraically independent
  \(\bar b_1, \dots, \bar b_n \in B_K\) with \(B_K\) module-finite over
  \(K[\bar b_1, \dots, \bar b_n]\). Clearing denominators, each \(\bar b_j\) is
  represented by some \(b_j \in B\) (writing \(1 \otimes b\) as \(b\)).

  \emph{Step 2 — clear the integral-dependence denominators to a single \(g\).} Fix a
  finite set \(\sigma\) of algebra generators of \(B_K\) over
  \(K[\bar b_1, \dots, \bar b_n]\). By Step 1 each generator's image
  \(1 \otimes x \in B_K\) is integral over \(K[\bar b_1, \dots, \bar b_n]\). For each
  \(x \in \sigma\), a standard integrality-and-localization descent argument yields a
  non-zero \(a_x \in A\) such that \(a_x \cdot (1 \otimes x)\)
  is integral over \(A[b_1, \dots, b_n]\). Multiplying over the finite
  set, set \(g_1 := \prod_{x \in \sigma} a_x\) and \(g := g_0 \cdot g_1\) (with \(g_0\) the
  Step-1 denominator witness). After inverting \(g\), every chosen generator becomes
  integral over \(A_g[b_1, \dots, b_n]\) and the generators generate \(B_g\) as an algebra;
  hence \(B_g\) is module-finite over \(A_g[b_1, \dots, b_n]\). (The single-coefficient
  lemma \cref{lem:gf_clear_one_denominator} provides the same denominator-clearing primitive
  at the level of one coefficient polynomial; the proof discharges a whole generator
  at once using this integrality-descent step, folding that primitive over the coefficients
  internally.)

  \emph{Step 3 — AlgHom assembly.} The algebra homomorphism \(\varphi\) of the statement
  is the \(A_g\)-algebra map \(A_g[X_1, \dots, X_n] \to B_g\) sending \(X_j \mapsto b_j\),
  where \(b_j \in B_g\) is the image of \(b_j \in B\). We establish injectivity and
  module-finiteness through a comparison with the situation over \(K\).

  \emph{Step 3a — the comparison map \(\nu : B_g \to B_K\).} Since \(g \ne 0\) in the
  domain \(A\) and \(A \subseteq K\), the image of \(g\) is a unit of the field \(K\), so
  it is invertible in \(B_K = K \otimes_A B\). Hence the canonical \(A\)-algebra map
  \(B \to B_K\), \(b \mapsto 1 \otimes b\), sends \(g\) to a unit and, by the universal
  property of localisation, factors through a canonical \(A\)-algebra map
  \[
    \nu : B_g = B[g^{-1}] \longrightarrow B_K, \qquad \nu(b/g^k) = (1 \otimes b)\,(1 \otimes g)^{-k}.
  \]
  The map \(\nu\) is injective. Indeed, \(B\) is a domain, and both \(B_g\) and \(B_K\) are
  localisations of \(B\) at multiplicative subsets of the non-zerodivisors: \(B_g\) inverts
  the powers of \(g\), while \(B_K = K \otimes_A B\) inverts (in particular) the images of
  all non-zero elements of \(A\), a multiplicative set containing the powers of \(g\). A
  localisation of a domain at a set of non-zerodivisors is injective into any further such
  localisation. Concretely, the natural map \(B \to B_K\) is injective on the domain \(B\)
  (its kernel is the set of \(b\) annihilated by some non-zero element of \(A\), which in a
  domain over which \(A\) acts faithfully is \(0\)), so \(\nu(b/g^k) = 0\) forces
  \(1 \otimes b = 0\), hence \(b = 0\) and \(b/g^k = 0\). Under \(\nu\) the lifts
  \(b_j \in B_g\) of the Noether generators map to the elements \(\bar b_j \in B_K\) of
  Step 1: \(\nu(b_j) = \bar b_j\).

  \emph{Step 3b — injectivity of \(\varphi\) by algebraic-independence descent
  \(K \to A_g\).} Injectivity of the \(A_g\)-algebra map
  \(\varphi : A_g[X_1, \dots, X_n] \to B_g\), \(X_j \mapsto b_j\), is equivalent to the
  algebraic independence of the family \((b_j)\) over \(A_g\): a non-zero element of the
  kernel is exactly a non-trivial \(A_g\)-polynomial relation \(P(b_1, \dots, b_n) = 0\).
  The localisation \(A_g = A[g^{-1}] \subseteq K\) embeds in the fraction field, giving an
  injection of \(A_g\)-algebras \(A_g \hookrightarrow K\). Now suppose
  \(P \in A_g[X_1, \dots, X_n]\) is non-zero with \(P(b_1, \dots, b_n) = 0\) in \(B_g\).
  Applying the injective map \(\nu\) and using \(\nu(b_j) = \bar b_j\) together with the fact
  that \(\nu\) is an \(A\)-algebra map extending scalars to \(K\), we obtain
  \(P^K(\bar b_1, \dots, \bar b_n) = 0\) in \(B_K\), where \(P^K\) is the image of \(P\)
  under the coefficient inclusion \(A_g[X_1, \dots, X_n] \hookrightarrow K[X_1, \dots, X_n]\).
  Since \(A_g \hookrightarrow K\) is injective, \(P \ne 0\) gives \(P^K \ne 0\); but the
  \(\bar b_j\) are algebraically independent over \(K\) (Step 1), so the only \(K\)-polynomial
  vanishing on them is \(0\), a contradiction. Hence no non-trivial relation exists, \((b_j)\)
  is algebraically independent over \(A_g\), and \(\varphi\) is injective.

  \emph{Step 3c — module-finiteness of \(B_g\) over \(\operatorname{im}\varphi\).} By Step 1,
  \(B_K\) is module-finite over \(K[\bar b_1, \dots, \bar b_n]\); fix a finite set of algebra
  generators of \(B_K\) over \(K[\bar b_1, \dots, \bar b_n]\) (Step 2). Each such generator
  satisfies a monic equation of integral dependence whose coefficients are polynomials in
  \(K[\bar b_1, \dots, \bar b_n]\). The element \(g\) of Step 2 was chosen as a common
  denominator for the (finitely many) coefficient polynomials of these equations, so that
  after inverting \(g\) every such coefficient already lies in
  \(A_g[b_1, \dots, b_n] = \operatorname{im}\varphi\). Each chosen generator is the
  \(\nu\)-image of an element of \(B_g\), and its integral-dependence equation is the
  \(\nu\)-image of the corresponding monic equation holding in \(B_g\) with coefficients in
  \(\operatorname{im}\varphi\); since \(\nu\) is injective, that equation holds in \(B_g\)
  itself. Thus each generator of \(B_g\) is integral over the subalgebra
  \(\operatorname{im}\varphi\), and these generators generate \(B_g\) as an
  \(\operatorname{im}\varphi\)-algebra. A finite-type algebra that is integral over a
  subalgebra is module-finite over it; hence \(B_g\) is module-finite over
  \(\operatorname{im}\varphi = A_g[b_1, \dots, b_n]\). Transported along the injection
  \(\varphi\), this is exactly module-finiteness of \(B_g\) over \(A_g[X_1, \dots, X_n]\) for
  the algebra structure that \(\varphi\) puts on \(B_g\), as claimed. This recovers the
  source's conclusion that \(B_g\) is finite over \(A_g[b_1, \dots, b_n]\).
\end{proof}

\begin{lemma}\leanok
[injectivity of a localisation lift]
  \label{lem:gf_isLocalization_lift_injective}
  \lean{AlgebraicGeometry.GenericFreeness.isLocalization_lift_injective}
  Let \(S\) be a localisation of a ring \(R\) at a multiplicative set \(M\), and let
  \(g : R \to P\) be a ring homomorphism sending every element of \(M\) to a unit of
  \(P\), so that it induces the lift \(\widehat g : S \to P\). If both the localisation
  map \(R \to S\) and \(g\) itself are injective, then \(\widehat g\) is injective.
\end{lemma}
\begin{proof}

\leanok
  The lift is characterised by \(\widehat g(r/m) = g(r)\,g(m)^{-1}\). By the universal
  criterion for injectivity of a localisation lift, \(\widehat g\) is injective iff for
  all \(r_1, r_2 \in R\) the equality of localisation images
  \((r_1 \mapsto S) = (r_2 \mapsto S)\) implies \(g(r_1) = g(r_2)\). Both the hypothesis
  and the conclusion of that implication reduce, via injectivity of \(R \to S\) and of
  \(g\), to the single condition \(r_1 = r_2\); hence the implication holds and
  \(\widehat g\) is injective. (This helper is applied twice in
  \cref{lem:gf_noether_clear_denominators}, to the comparison maps \(\nu : B_g \to B_K\)
  and \(\psi : A_g \to K\).)
\end{proof}

\begin{lemma}


\leanok
[L5a — the generic-rank short exact sequence]
  \label{lem:gf_generic_rank_ses}
  \lean{AlgebraicGeometry.GenericFreeness.gf_generic_rank_ses}
  % SOURCE: [Nitsure], \S4 ``Lemma on Generic Flatness'', proof, inductive step
  % (read from references/nitsure-hilbert-quot-src/nitsure-hilbert-quot.tex, L1760-L1765).
  % SOURCE QUOTE PROOF:
  % "Let $m$ be the generic rank of the finite module $B_g$ over the
  %  domain $A_g[b_1,\ldots, b_n]$. Then we have a short exact
  %  sequence of $A_g[b_1,\ldots, b_n]$-modules of the form
  %  $$0\to A_g[b_1,\ldots, b_n]^{\oplus m} \to B_g \to T \to 0$$
  %  where $T$ is a finite torsion module over $A_g[b_1,\ldots, b_n]$."
  \textit{Source: [Nitsure], \S4, ``Lemma on Generic Flatness'' (generic-rank SES).}
  Let \(A\) be a noetherian domain, \(d \ge 0\), and \(N\) a finite module over the
  polynomial ring \(P_d := A[X_1, \dots, X_d]\) (with its induced \(A\)-module structure
  through the scalar tower \(A \to P_d \to N\)). Then there exist \(m \in \mathbb{N}\) and
  an injective \(P_d\)-linear map \(\varphi : P_d^{\oplus m} \to N\) whose cokernel
  \(T := N / \operatorname{im}\varphi\) is a torsion \(P_d\)-module. Equivalently, there is
  a short exact sequence of \(P_d\)-modules
  \[
    0 \to P_d^{\oplus m} \xrightarrow{\ \varphi\ } N \to T \to 0,
    \qquad T \text{ torsion}.
  \]

  The integer \(m\) is the \emph{generic rank} of \(N\): the dimension of the
  \(\mathrm{Frac}(P_d)\)-vector space
  \(\mathrm{Frac}(P_d) \otimes_{P_d} N\), i.e.\ of the localisation of \(N\) at the
  non-zero-divisors of the domain \(P_d\). This step lives over \(P_d\) directly —
  \emph{no inversion of any element \(g \in A\) is required here}; the element \(g\) of
  Nitsure's argument enters only at the reindexing step \cref{lem:gf_torsion_reindex}.
\end{lemma}

\begin{proof}

\leanok
  % SOURCE QUOTE PROOF: see the comment above the statement (L1760-L1765).
  \textit{Source: [Nitsure], \S4 (generic-rank SES).}
  Write \(P_d = A[X_1, \dots, X_d]\); since \(A\) is a noetherian domain, so is \(P_d\)
  (Hilbert basis theorem, and a polynomial ring over a domain is a domain). Let
  \(Q = \mathrm{Frac}(P_d)\) be its fraction field and
  \(N_Q := Q \otimes_{P_d} N\) the localisation of \(N\) at the non-zero-divisors of
  \(P_d\); as \(N\) is finite over \(P_d\), \(N_Q\) is a finite-dimensional
  \(Q\)-vector space. Put \(m := \dim_Q N_Q\), the generic rank.

  Choose a \(Q\)-basis of \(N_Q\) and clear denominators: each basis vector is a
  \(Q\)-unit multiple of the image of some \(v_i \in N\), so the images of
  \(v_1, \dots, v_m \in N\) again form a \(Q\)-basis of \(N_Q\). In particular the family
  \((v_i)_{i}\) is linearly independent over \(Q\); restricting scalars along the
  injection \(P_d \hookrightarrow Q\), it is linearly independent over \(P_d\). Let
  \(\varphi : P_d^{\oplus m} \to N\) be the \(P_d\)-linear map
  \((c_1, \dots, c_m) \mapsto \sum_i c_i v_i\). Linear independence of \((v_i)\) over
  \(P_d\) makes \(\varphi\) injective.

  Set \(T := N / \operatorname{im}\varphi\). Because the images of \(v_1, \dots, v_m\)
  already span \(N_Q\) over \(Q\), the localised map \(\varphi \otimes Q : Q^{\oplus m}
  \to N_Q\) is an isomorphism, so \(Q \otimes_{P_d} T = 0\). A finite \(P_d\)-module whose
  localisation at the non-zero-divisors of \(P_d\) vanishes is torsion; hence \(T\) is a
  torsion \(P_d\)-module. This is the asserted short exact sequence.
\end{proof}

\begin{lemma}[L5b anchor — a finite torsion module has a non-zero-divisor annihilator]
  \label{lem:annihilator_meets_nonZeroDivisors}
  \lean{Submodule.annihilator_top_inter_nonZeroDivisors}
  \mathlibok
  % SOURCE: [Nitsure], \S4 ``Lemma on Generic Flatness'', proof, base case $n=-1$
  % (read from references/nitsure-hilbert-quot-src/nitsure-hilbert-quot.tex, L1724-L1727)
  % -- the same torsion-annihilation primitive, here provided in general by Mathlib.
  % SOURCE QUOTE:
  % "each $v\in M$ is annihilated by some non-zero element of $A$. Taking a finite
  %  generating set, and a common multiple of corresponding annihilating
  %  elements, we see there exists an $f\ne 0$ in $A$ with $fM=0$."
  \textit{Provided by Mathlib.}
  Let \(R\) be a commutative ring and \(M\) a finite torsion \(R\)-module. Then the
  annihilator of \(M\) meets the non-zero-divisors of \(R\): the set
  \((\top : \operatorname{Submodule} R\,M).\mathrm{annihilator} \cap R^{\circ}\) is
  non-empty, where \(R^{\circ}\) denotes the non-zero-divisors of \(R\).
\end{lemma}

\begin{lemma}\leanok
[L5b.1 — annihilator extraction for the torsion module]
  \label{lem:gf_torsion_annihilator}
  \lean{AlgebraicGeometry.GenericFreeness.gf_torsion_annihilator}
  \uses{lem:annihilator_meets_nonZeroDivisors}
  % SOURCE: [Nitsure], \S4 ``Lemma on Generic Flatness'', proof, base case $n=-1$
  % (read from references/nitsure-hilbert-quot-src/nitsure-hilbert-quot.tex, L1724-L1727).
  % SOURCE QUOTE PROOF:
  % "each $v\in M$ is annihilated by some non-zero element of $A$. Taking a finite
  %  generating set, and a common multiple of corresponding annihilating
  %  elements, we see there exists an $f\ne 0$ in $A$ with $fM=0$."
  \textit{Source: [Nitsure], \S4, ``Lemma on Generic Flatness'' (torsion annihilation).}
  Let \(A\) be a noetherian domain, \(d \ge 0\), and \(T\) a finite \emph{torsion} module
  over \(P_d := A[X_1, \dots, X_d]\). Then there exists a non-zero
  \(F \in \operatorname{Ann}_{P_d}(T)\); consequently the \(P_d\)-action on \(T\) factors
  through \(P_d/(F)\) and \(T\) is a finite module over \(P_d/(F)\).
\end{lemma}

\begin{proof}

\leanok
  % SOURCE QUOTE PROOF: see the comment above the statement (L1724-L1727).
  \textit{Source: [Nitsure], \S4 (torsion annihilation).}
  The ring \(P_d = A[X_1, \dots, X_d]\) is a noetherian domain (a polynomial ring over the
  noetherian domain \(A\)), and \(T\) is a finite torsion \(P_d\)-module. By
  \cref{lem:annihilator_meets_nonZeroDivisors} the annihilator of \(T\) meets the
  non-zero-divisors of \(P_d\); choose \(F\) in that intersection. In the domain \(P_d\) a
  non-zero-divisor is exactly a non-zero element, so \(F \ne 0\), and \(F \cdot T = 0\).
  Since \(F\) annihilates \(T\), the \(P_d\)-module structure descends along the quotient
  map \(P_d \twoheadrightarrow P_d/(F)\); the image of any finite \(P_d\)-generating set of
  \(T\) generates \(T\) over \(P_d/(F)\), so \(T\) is finite over \(P_d/(F)\).
\end{proof}

\subsubsection{Nagata change-of-variables machinery}
\label{sec:gf_nagata_machinery}

The Nagata normalisation lemma \cref{lem:gf_nagata_monic_lastVar} is assembled from
technical helper lemmas implementing the triangular change of variables and the
radix-encoding degree bookkeeping, following the Nagata trick of
\cite{nitsure-hilbert-quot}~\S4 over a noetherian domain. Throughout, \(k\) is a domain,
\(n \ge 0\), and \(f \in k[X_0, \dots, X_n]\) is a fixed polynomial; we write
\(\mathrm{up} := 2 + \deg_{\mathrm{tot}}(f)\) and the radix weights \(r_i := \mathrm{up}^{\,i}\)
for \(i = 0, \dots, n\) (so every exponent occurring in \(f\) is a base-\(\mathrm{up}\) digit).
The distinguished variable is \(X_0\).

\begin{definition}\leanok
[Nagata substitution \(T_1\)]
  \label{def:gf_nagata_T1}
  \lean{AlgebraicGeometry.GenericFreeness.T1}
  \textit{Source: Nitsure \S4 (Nagata change of variables).}
  For a scalar \(c \in k\), the substitution \(T_1(c)\) is the \(k\)-algebra endomorphism
  of \(k[X_0, \dots, X_n]\) determined on generators by
  \[
    X_0 \mapsto X_0, \qquad X_i \mapsto X_i + c\,X_0^{\,r_i} \quad (i \ne 0).
  \]
  It is triangular in the variables \(X_i\) (\(i \ne 0\)) over the distinguished variable
  \(X_0\); this is the device realising the support-dimension drop.
\end{definition}

\begin{lemma}\leanok
[\(T_1\) is invertible by the opposite shift]
  \label{lem:gf_t1_comp_t1_neg}
  \lean{AlgebraicGeometry.GenericFreeness.t1_comp_t1_neg}
  \uses{def:gf_nagata_T1}
  For every \(c \in k\), \(T_1(c) \circ T_1(-c) = \mathrm{id}\); the Nagata substitution is
  invertible, with inverse the substitution of opposite coefficient.
\end{lemma}
\begin{proof}
  \uses{def:gf_nagata_T1}
  \leanok
  It suffices to check the two composites agree on each generator: \(X_0 \mapsto X_0\), and
  \(X_i \mapsto X_i + (-c)X_0^{r_i} + c\,X_0^{r_i} = X_i\) for \(i \ne 0\).
\end{proof}

\begin{definition}\leanok
[Nagata automorphism \(T\)]
  \label{def:gf_nagata_T}
  \lean{AlgebraicGeometry.GenericFreeness.T}
  \uses{def:gf_nagata_T1, lem:gf_t1_comp_t1_neg}
  \textit{Source: Nitsure \S4 (Nagata change of variables).}
  The Nagata change of variables is the \(k\)-algebra automorphism \(T := T_1(1)\) of
  \(k[X_0, \dots, X_n]\) — i.e.\ \(X_i \mapsto X_i + X_0^{\,r_i}\) (\(i \ne 0\)) and
  \(X_0 \mapsto X_0\) — packaged as an algebra equivalence with inverse \(T_1(-1)\) via
  \cref{lem:gf_t1_comp_t1_neg}.
\end{definition}

\begin{lemma}\leanok
[Exponent bound is preserved by listing]
  \label{lem:gf_lt_up}
  \lean{AlgebraicGeometry.GenericFreeness.lt_up}
  If every entry of an exponent vector \(v\) satisfies \(v_i < \mathrm{up}\), then
  every entry of the list of values of \(v\) is also \(< \mathrm{up}\).
\end{lemma}
\begin{proof}

\leanok
  The list of values of \(v\) consists exactly of the entries \(v_i\); each satisfies
  \(v_i < \mathrm{up}\) by hypothesis.
\end{proof}

\begin{lemma}\leanok
[Injectivity of the radix encoding]
  \label{lem:gf_sum_r_mul_ne}
  \lean{AlgebraicGeometry.GenericFreeness.sum_r_mul_ne}
  \uses{lem:gf_lt_up}
  Let \(v, w\) be exponent vectors with all entries \(< \mathrm{up}\) and \(v \ne w\). Then
  the weighted base-\(\mathrm{up}\) sums differ:
  \(\sum_i r_i v_i \ne \sum_i r_i w_i\), where \(r_i = \mathrm{up}^{\,i}\).
\end{lemma}
\begin{proof}
  \uses{lem:gf_lt_up}
  \leanok
  The map \(v \mapsto \sum_i \mathrm{up}^{\,i} v_i\) is the value of the base-\(\mathrm{up}\)
  digit string \(v\); since every digit is \(< \mathrm{up}\) (\cref{lem:gf_lt_up}), this
  encoding is injective by uniqueness of base-\(\mathrm{up}\) expansions, so distinct
  \(v \ne w\) give distinct sums.
\end{proof}

\begin{lemma}\leanok
[\(X_0\)-degree of a transformed monomial]
  \label{lem:gf_degreeOf_zero_t}
  \lean{AlgebraicGeometry.GenericFreeness.degreeOf_zero_t}
  \uses{def:gf_nagata_T}
  For a non-zero scalar \(a \ne 0\) and exponent vector \(v\), the degree in \(X_0\) of the
  transformed monomial is the radix value
  \[
    \deg_{X_0}\!\big(T(a\,X^v)\big) = \sum_i r_i v_i, \qquad r_i = \mathrm{up}^{\,i}.
  \]
\end{lemma}
\begin{proof}
  \uses{def:gf_nagata_T}
  \leanok
  Under \(T\), \(a\,X^v \mapsto a\,X_0^{v_0}\prod_{i \ne 0}(X_i + X_0^{r_i})^{v_i}\). Read as
  a polynomial in \(X_0\) over \(k[X_1, \dots, X_n]\), each factor
  \((X_i + X_0^{r_i})^{v_i}\) has \(X_0\)-degree \(r_i v_i\) and leading coefficient \(1\);
  the degrees add to \(\sum_i r_i v_i\) (the hypothesis \(a \ne 0\) keeps the top term from
  vanishing over the domain \(k\)).
\end{proof}

\begin{lemma}\leanok
[Distinct monomials acquire distinct \(X_0\)-degrees]
  \label{lem:gf_degreeOf_t_ne_of_ne}
  \lean{AlgebraicGeometry.GenericFreeness.degreeOf_t_ne_of_ne}
  \uses{lem:gf_degreeOf_zero_t, lem:gf_sum_r_mul_ne}
  For distinct support exponents \(v \ne w\) of \(f\), the transformed monomials
  \(T(c_v X^v)\) and \(T(c_w X^w)\) (where \(c_v, c_w\) are the corresponding coefficients
  of \(f\)) have distinct degrees in \(X_0\), making the top \(X_0\)-degree of \(T(f)\) unique.
\end{lemma}
\begin{proof}
  \uses{lem:gf_degreeOf_zero_t, lem:gf_sum_r_mul_ne}
  \leanok
  By \cref{lem:gf_degreeOf_zero_t} the two degrees are \(\sum_i r_i v_i\) and
  \(\sum_i r_i w_i\). Each support exponent is bounded by the total degree of \(f\), hence
  \(< \mathrm{up} = 2 + \deg_{\mathrm{tot}} f\); radix-injectivity
  \cref{lem:gf_sum_r_mul_ne} then yields \(\sum_i r_i v_i \ne \sum_i r_i w_i\).
\end{proof}

\begin{lemma}\leanok
[Leading coefficient of a transformed monomial]
  \label{lem:gf_leadingCoeff_finSuccEquiv_t}
  \lean{AlgebraicGeometry.GenericFreeness.leadingCoeff_finSuccEquiv_t}
  \uses{def:gf_nagata_T, def:gf_nagata_T1}
  For any exponent \(v\) with coefficient \(c_v := [X^v]\,f\), viewing \(T(c_v X^v)\) as a
  univariate polynomial in \(X_0\) over \(k[X_1, \dots, X_n]\), its leading coefficient in
  \(X_0\) is the constant \(c_v \in k\).
\end{lemma}
\begin{proof}
  \uses{def:gf_nagata_T, def:gf_nagata_T1}
  \leanok
  Each factor \((X_i + X_0^{r_i})^{v_i}\) is monic in \(X_0\) with leading coefficient
  \(1\); a product of monics is monic, so the leading coefficient of \(T(c_v X^v)\) in
  \(X_0\) is exactly the scalar \(c_v\).
\end{proof}

\begin{lemma}\leanok
[Leading coefficient of \(T(f)\) is a non-zero constant]
  \label{lem:gf_T_leadingcoeff_eq}
  \lean{AlgebraicGeometry.GenericFreeness.T_leadingcoeff_eq}
  \uses{lem:gf_degreeOf_t_ne_of_ne, lem:gf_leadingCoeff_finSuccEquiv_t}
  Suppose \(f \ne 0\). Then there is a support exponent \(v \in \mathrm{supp}\,f\) such that
  the leading coefficient of \(T(f)\) in \(X_0\) equals the constant \(c_v := [X^v]\,f \ne 0\).
  Thus \(T(f)\) is, up to the non-zero scalar \(c_v\), monic in \(X_0\).
\end{lemma}
\begin{proof}
  \uses{lem:gf_degreeOf_t_ne_of_ne, lem:gf_leadingCoeff_finSuccEquiv_t}
  \leanok
  Expand \(f = \sum_{w \in \mathrm{supp}\,f} c_w X^w\). By \cref{lem:gf_degreeOf_t_ne_of_ne}
  the transformed monomials \(T(c_w X^w)\) have pairwise distinct \(X_0\)-degrees, so one of
  them — say at \(v\) — strictly dominates the others; the leading term of
  \(T(f) = \sum_w T(c_w X^w)\) is then that of \(T(c_v X^v)\), whose leading coefficient is
  the constant \(c_v\) by \cref{lem:gf_leadingCoeff_finSuccEquiv_t}. Since
  \(v \in \mathrm{supp}\,f\), \(c_v \ne 0\).
\end{proof}

\begin{lemma}\leanok
[Naturality of the variable-splitting isomorphism under base change]
  \label{lem:gf_finSuccEquiv_map_comm}
  \lean{AlgebraicGeometry.GenericFreeness.finSuccEquiv_map_comm}
  For any ring homomorphism \(\varphi : A \to B\) and polynomial \(q \in A[X_0, \dots, X_m]\),
  the isomorphism \(A[X_0, \dots, X_m] \cong (A[X_1, \dots, X_m])[X_0]\) is natural with
  respect to applying \(\varphi\) coefficientwise.
\end{lemma}
\begin{proof}

\leanok
  By induction on \(q\): the isomorphism is additive and multiplicative, and agrees with
  \(\varphi\) on constants and variables, so it commutes with \(\varphi\) on all polynomials.
\end{proof}

\begin{lemma}\leanok
[Constant inclusion under the variable-splitting isomorphism]
  \label{lem:gf_finSuccEquiv_rename_succ}
  \lean{AlgebraicGeometry.GenericFreeness.finSuccEquiv_rename_succ}
  For \(s \in R[X_1, \dots, X_n]\), the composite of the inclusion
  \(R[X_1, \dots, X_n] \hookrightarrow R[X_0, \dots, X_n]\) (sending \(X_j \mapsto X_j\))
  with the isomorphism \(R[X_0, \dots, X_n] \cong (R[X_1, \dots, X_n])[X_0]\) sends \(s\)
  to the constant polynomial \(s \in (R[X_1, \dots, X_n])[X_0]\).
\end{lemma}
\begin{proof}

\leanok
  By induction on \(s\): a constant and each variable \(X_j\) (\(j \ge 1\)) map to themselves,
  and the isomorphism is an algebra map, so it sends products and sums of such to constants in
  the outer polynomial ring.
\end{proof}

\begin{lemma}\leanok
[L5b.2 — Nagata change of variables: monic in the distinguished variable]
  \label{lem:gf_nagata_monic_lastVar}
  \lean{AlgebraicGeometry.GenericFreeness.gf_nagata_monic_lastVar}
  \uses{def:gf_nagata_T, lem:gf_T_leadingcoeff_eq, lem:gf_finSuccEquiv_map_comm}
  % SOURCE: [Nitsure], \S4 ``Lemma on Generic Flatness'', proof, inductive step
  % (read from references/nitsure-hilbert-quot-src/nitsure-hilbert-quot.tex, L1766-L1768);
  % the Nagata substitution is the standard device realising the support-dimension drop.
  % SOURCE QUOTE PROOF:
  % "Therefore, the dimension of the support of $K\otimes_{A_g}T$
  %  as a $K\otimes_{A_g}(B_g)$-module is strictly less than $n$. Hence by
  %  induction on $n$ (applied to the data $A_g$, $B_g$, $T$),"
  \textit{Source: [Nitsure], \S4, ``Lemma on Generic Flatness'' (support-dimension drop).}
  Let \(A\) be a domain, \(n \ge 0\), and \(0 \ne F \in A[X_0, \dots, X_n]\). Then there
  exist a triangular \(A\)-algebra automorphism \(e\) of \(A[X_0, \dots, X_n]\) and a
  non-zero \(g \in A\) such that, after inverting \(g\), the image of \(e(F)\) in
  \(A_g[X_0, \dots, X_n]\), viewed as a univariate polynomial in \(X_0\) over
  \(A_g[X_1, \dots, X_n]\), has a \emph{unit} leading coefficient. Equivalently, \(e(F)\)
  is monic in \(X_0\) up to the unit \(g\) over \(A_g\).
\end{lemma}

\begin{proof}

\leanok
  % SOURCE QUOTE PROOF: see the comment above the statement (L1766-L1768).
  \textit{Source: [Nitsure], \S4 (support-dimension drop, Nagata normalisation).}
  Apply the Nagata substitution \(X_i \mapsto X_i + X_0^{r^{\,i}}\) for \(i = 1, \dots, n\)
  and \(X_0 \mapsto X_0\), with \(r\) larger than every exponent occurring in \(F\): it is a
  triangular (hence invertible) \(A\)-algebra automorphism \(e\) of \(A[X_0, \dots, X_n]\).
  Under this substitution the distinct monomials of \(F\) acquire distinct degrees in
  \(X_0\), so \(e(F)\), expanded as a polynomial in \(X_0\), has a unique top-degree term
  whose coefficient is the coefficient \(c \in A\) of the top total-degree monomial of
  \(F\); since \(F \ne 0\) this \(c \ne 0\). Put \(g := c\). The degree of \(e(F)\) in
  \(X_0\) equals the total degree of \(F\), and its leading coefficient in \(X_0\) is the
  constant \(g\). Passing to \(A_g = A[g^{-1}]\), the image of \(g\) is a unit, so the
  image of \(e(F)\) in \(A_g[X_0, \dots, X_n]\) is monic in \(X_0\) up to the unit \(g\).
\end{proof}

\begin{lemma}[L5b engine anchor — quotient by a monic univariate polynomial is finite]
  \label{lem:polynomial_monic_quotient_finite}
  \lean{Polynomial.Monic.finite_quotient}
  \mathlibok
  \textit{Provided by Mathlib.}
  Let \(S\) be a commutative ring and \(f \in S[X]\) a monic polynomial. Then the quotient
  \(S[X]/(f)\) is a module-finite \(S\)-module (indeed free of rank \(\deg f\)).
\end{lemma}

\begin{lemma}\leanok
[L5b.3 — single-variable elimination engine (shared)]
  \label{lem:gf_mvPolynomial_quotient_finite_monic}
  \lean{AlgebraicGeometry.GenericFreeness.mvPolynomial_quotient_finite_of_monic_lastVar}
  \uses{lem:polynomial_monic_quotient_finite, lem:gf_finSuccEquiv_rename_succ}
  % SOURCE: [Nitsure], \S4 ``Lemma on Generic Flatness'', proof, inductive step
  % (read from references/nitsure-hilbert-quot-src/nitsure-hilbert-quot.tex, L1766-L1768);
  % the division algorithm in one variable is the standard mechanism of the support drop.
  % SOURCE QUOTE PROOF:
  % "Therefore, the dimension of the support of $K\otimes_{A_g}T$
  %  as a $K\otimes_{A_g}(B_g)$-module is strictly less than $n$. Hence by
  %  induction on $n$ (applied to the data $A_g$, $B_g$, $T$),"
  \textit{Source: [Nitsure], \S4, ``Lemma on Generic Flatness'' (single-variable elimination).}
  Let \(R\) be a commutative ring, \(n \in \mathbb{N}\), and \(p \in R[X_0, \dots, X_n]\) a
  polynomial whose image under the isomorphism
  \(R[X_0, \dots, X_n] \cong (R[X_1, \dots, X_n])[X_0]\) has a \emph{unit} leading
  coefficient — i.e.\ \(p\) is monic in \(X_0\) up to a unit. Then the
  quotient \(R[X_0, \dots, X_n]/(p)\) is module-finite over the coefficient subring
  \(R[X_1, \dots, X_n]\) (embedded by \(X_j \mapsto X_j\), \(j = 1, \dots, n\)).
\end{lemma}

\begin{proof}

\leanok
  % SOURCE QUOTE PROOF: see the comment above the statement (L1766-L1768).
  \textit{Source: [Nitsure], \S4 (single-variable elimination).}
  Write \(S := R[X_1, \dots, X_n]\) and let
  \(\Phi : R[X_0, \dots, X_n] \xrightarrow{\ \sim\ } S[X_0]\) be the standard \(R\)-algebra
  isomorphism, which sends \(X_0\) to the univariate variable and embeds \(S\) as the constants. Put \(q := \Phi(p) \in S[X_0]\); by hypothesis its leading
  coefficient \(u\) is a unit, so the associate \(\tilde q := u^{-1} \cdot q\) is monic and
  \((\tilde q) = (q)\) in \(S[X_0]\). The isomorphism \(\Phi\) carries the ideal \((p)\) onto
  \((q) = (\tilde q)\), hence descends to a ring isomorphism
  \(R[X_0, \dots, X_n]/(p) \cong S[X_0]/(\tilde q)\) that is \(S\)-linear, because \(S\) sits
  inside both sides as the constants (respectively as the coefficient subring
  \(R[X_1, \dots, X_n]\)). By \cref{lem:polynomial_monic_quotient_finite} the quotient
  \(S[X_0]/(\tilde q)\) is module-finite over \(S\); transporting finiteness along the
  \(S\)-linear isomorphism shows \(R[X_0, \dots, X_n]/(p)\) is module-finite over
  \(S = R[X_1, \dots, X_n]\).
\end{proof}

\begin{lemma}


\leanok
[L5b — torsion reindex onto fewer variables]
  \label{lem:gf_torsion_reindex}
  \lean{AlgebraicGeometry.GenericFreeness.gf_torsion_reindex}
  \uses{lem:gf_torsion_annihilator, lem:gf_nagata_monic_lastVar,
        lem:gf_mvPolynomial_quotient_finite_monic, lem:gf_pullback_module_transport,
        lem:gf_finite_of_quotient_ringequiv, lem:gf_islocalizedmodule_restrictscalars}
  % SOURCE: [Nitsure], \S4 ``Lemma on Generic Flatness'', proof, inductive step
  % (read from references/nitsure-hilbert-quot-src/nitsure-hilbert-quot.tex, L1766-L1768).
  % SOURCE QUOTE PROOF:
  % "Therefore, the dimension of the support of $K\otimes_{A_g}T$
  %  as a $K\otimes_{A_g}(B_g)$-module is strictly less than $n$. Hence by
  %  induction on $n$ (applied to the data $A_g$, $B_g$, $T$),"
  \textit{Source: [Nitsure], \S4, ``Lemma on Generic Flatness'' (support-dimension drop).}
  Let \(A\) be a noetherian domain, \(d \ge 1\), and \(T\) a finite \emph{torsion} module
  over \(P_d := A[X_1, \dots, X_d]\) (with its induced \(A\)-module structure). Then there
  exist a non-zero \(g \in A\) and an integer \(m' < d\) such that, after inverting \(g\),
  the localised module \(T_g := T[g^{-1}]\) is module-finite over the polynomial ring
  \(A_g[X_1, \dots, X_{m'}]\) in \(m'\) variables over the base ring \(A_g = A[g^{-1}]\);
  that is, \(T_g\) carries compatible \(A_g\)- and \(A_g[X_1, \dots, X_{m'}]\)-module
  structures (via the scalar tower \(A_g \to A_g[X_1, \dots, X_{m'}]\)) and is finite over
  \(A_g[X_1, \dots, X_{m'}]\). One may always take \(m' = d - 1\).
\end{lemma}

\begin{proof}

\leanok
  % SOURCE QUOTE PROOF: see the comment above the statement (L1766-L1768).
  \textit{Source: [Nitsure], \S4 (support-dimension drop).}
  \uses{lem:gf_torsion_annihilator, lem:gf_nagata_monic_lastVar,
        lem:gf_mvPolynomial_quotient_finite_monic, lem:gf_pullback_module_transport,
        lem:gf_finite_of_quotient_ringequiv, lem:gf_islocalizedmodule_restrictscalars}
  Write \(d = (d-1) + 1\) and \(P_d = A[X_0, \dots, X_{d-1}]\) (the distinguished variable
  \(X_0\) is the one eliminated in the single-variable step; cf.\ \cref{lem:gf_nagata_monic_lastVar}).

  \emph{Annihilator.} By \cref{lem:gf_torsion_annihilator} there is a non-zero
  \(F \in \operatorname{Ann}_{P_d}(T)\), and \(T\) is a finite module over \(P_d/(F)\).

  \emph{Nagata normalisation.} Applying \cref{lem:gf_nagata_monic_lastVar} to \(F \ne 0\)
  yields a triangular \(A\)-algebra automorphism \(e\) of \(P_d\) and a non-zero
  \(g \in A\) such that the image of \(e(F)\) in \(A_g[X_0, \dots, X_{d-1}]\) has a unit
  leading coefficient in \(X_0\). The automorphism \(e\) induces a ring isomorphism
  \(P_d/(F) \cong P_d/(e\,F)\), so \(T\) is also finite over \(P_d/(e\,F)\). Localising at
  the powers of \(g\) (localisation commutes with the quotient), \(T_g\) is finite over
  \(P_d[g^{-1}]/(e\,F) = A_g[X_0, \dots, X_{d-1}]/(e\,F)\).

  \emph{Elimination.} Apply the shared engine
  \cref{lem:gf_mvPolynomial_quotient_finite_monic} with \(R = A_g\) and \(p\) the image of
  \(e(F)\) (whose leading coefficient in \(X_0\) is a unit): the quotient
  \(A_g[X_0, \dots, X_{d-1}]/(e\,F)\) is module-finite over the coefficient subring
  \(A_g[X_1, \dots, X_{d-1}]\), a polynomial ring in \(m' := d - 1\) variables over
  \(A_g\).

  \emph{Localisation transport.} The elimination step records module-finiteness of the
  \emph{quotient ring} over its coefficient subring, but the conclusion of the lemma is
  about a localisation of the torsion module \(T\). Making the two agree is the
  mathematical content of this step, which proceeds as follows.

  Two distinct localisations of \(T\) are in play and must be identified. The natural
  setting for the elimination is the localisation \(T'_g := T[\,\mathcal{C}^{-1}]\) of \(T\)
  at the multiplicative set \(\mathcal{C} = \{\,c \cdot g^{\,k} : k \ge 0\,\}\) of constant
  polynomials \(g^k\) inside \(P_d\); this localisation carries, for free, the full
  module structure over the localised polynomial ring
  \(P_d[g^{-1}] = A_g[X_0, \dots, X_{d-1}]\), its finiteness, and the quotient structure by
  \((e\,F)\). The \emph{goal}, however, is phrased with \(T_g := T[g^{-1}]\), the
  localisation of \(T\) at the powers of the single element \(g \in A\). One first proves
  \(T'_g\) module-finite over \(R := A_g[X_1, \dots, X_{m'}]\) and then transports that
  finiteness to \(T_g\).

  The change of variables \(e\) is transported to the localised polynomial ring as a
  \emph{ring isomorphism} \(\bar e : P_d[g^{-1}] \xrightarrow{\sim} P_d[g^{-1}]\), obtained
  by localising \(e\) at \(\mathcal{C}\) (which \(e\) preserves, since \(e\) fixes the
  constants \(c\,X^0\)); it carries the localised annihilator \(F_g\) to the
  Nagata-normalised generator \(G\) and fixes the constants coming from the base. Passing
  to quotients, \(\bar e\) descends to a ring isomorphism
  \(\psi : P_d[g^{-1}]/(F_g) \xrightarrow{\sim} P_d[g^{-1}]/(G)\) of the two quotient rings
  (using that the image of the ideal \((F_g)\) under \(\bar e\) is \((G)\)). Both quotient
  rings act on \(T'_g\), and module-finiteness over \(R\) of the elimination-side quotient
  \(P_d[g^{-1}]/(G)\) is transported across \(\psi\) — a ring isomorphism of the
  \emph{acting} ring compatible with the fixed \(R\)-algebra structure — by
  \cref{lem:gf_finite_of_quotient_ringequiv}, yielding \(T'_g\) module-finite over \(R\).

  Finally the two localisations are identified. Restricting scalars of the canonical
  \(P_d\)-localisation map of \(T\) along \(A \to P_d\) exhibits \(T'_g\) as a localisation
  of \(T\) \emph{at the powers of \(g\) over \(A\)} — this descent of the
  localised-module property along the scalar tower is
  \cref{lem:gf_islocalizedmodule_restrictscalars} — so the universal property produces a
  linear isomorphism \(T_g \cong T'_g\) identifying the two.

  \emph{Module structure and scalar tower.} The conclusion requires \(T_g\) to carry an
  \(A_g\)-module structure compatible with the scalar tower
  \(A_g \to A_g[X_1, \dots, X_{m'}] \to T_g\). The localisation \(T_g = T[g^{-1}]\) already
  carries the canonical \(A_g\)-module structure induced by localisation, and this is the
  structure required by the scalar tower. The identification \(T_g \cong T'_g\) must
  therefore be \(A_g\)-linear: the constant-polynomial action of \(A_g\) on \(T'_g\) factors
  through the original \(A\)-action on \(T\) in exactly the same way as the canonical
  \(g\)-localisation action does, so the two structures agree and the isomorphism is
  \(A_g\)-linear. The \(R\)-module structure and finiteness then pull back from \(T'_g\) to
  \(T_g\) via \cref{lem:gf_pullback_module_transport}, and the required scalar tower
  \(A_g \to R \to T_g\) follows.

  \emph{Transitivity.} Composing the two module-finite extensions, \(T_g\) is finite over
  \(A_g[X_1, \dots, X_{m'}]\). Since \(m' = d - 1 < d\) (using \(d \ge 1\)), this is the
  required reindexing; no Krull-dimension theory is needed, the number of variables
  dropping by exactly one through the single-variable elimination.
\end{proof}

\subsubsection{Module-structure transport helpers (reindex)}
\label{sec:gf_reindex_transport_helpers}

The localisation-transport step of \cref{lem:gf_torsion_reindex} moves a module structure
and its finiteness across (i) an additive equivalence of the underlying groups, (ii) a ring
isomorphism of the acting ring, and (iii) a change of localising monoid along a scalar
tower. The following helper lemmas package those three transport principles. They are general
module-theory facts; no external source is cited.

\begin{lemma}\leanok
[Module transport along an additive equivalence]
  \label{lem:gf_pullback_module_transport}
  \lean{AlgebraicGeometry.GenericFreeness.pullbackModuleAddEquiv}
  \lean{AlgebraicGeometry.GenericFreeness.finite_of_pullbackModuleAddEquiv}
  \lean{AlgebraicGeometry.GenericFreeness.pullback_isScalarTower}
  Let \(e : M \xrightarrow{\sim} N\) be an additive equivalence and let \(M\) carry an
  \(R\)-module structure. Pulling the action back along \(e\) by
  \(r \cdot y := e\big(r \cdot e^{-1}(y)\big)\) makes \(N\) an \(R\)-module for which \(e\)
  is \(R\)-linear. Under this transported structure: (a) if \(M\) is module-finite over
  \(R\) then so is \(N\); and (b) if \(M\) additionally carries a compatible action of a
  base ring \(A_g\) with scalar tower \(A_g \to R \to M\), then the pulled-back actions on
  \(N\) again form a scalar tower \(A_g \to R \to N\).
\end{lemma}
\begin{proof}

\leanok
  The module axioms for \(y \mapsto e(r \cdot e^{-1}(y))\) follow termwise from those on
  \(M\) together with the fact that \(e\) and \(e^{-1}\) are additive and mutually inverse;
  by construction \(e\) is then an \(R\)-linear isomorphism. Statement (a) follows by
  transporting module-finiteness across the \(R\)-linear isomorphism \(e\).
  Statement (b) checks the scalar-tower identity
  \((a\,r)\cdot y = a\cdot(r\cdot y)\) on \(N\) by transporting it across \(e\) and applying
  associativity of the scalar action on \(M\).
\end{proof}

\begin{lemma}\leanok
[Finite transport along a ring isomorphism of the acting ring]
  \label{lem:gf_finite_of_quotient_ringequiv}
  \lean{AlgebraicGeometry.GenericFreeness.finite_of_quotientRingEquiv}
  Let \(R\) be a commutative ring and let \(\psi : B_1 \xrightarrow{\sim} B_2\) be a ring
  isomorphism between two \(R\)-algebras that is compatible with the structure maps, i.e.\
  \(\psi \circ (\text{algebra map } R \to B_1) = (\text{algebra map } R \to B_2)\). If
  \(B_2\) is module-finite over \(R\) and \(M\) is module-finite over \(B_1\) (with the
  scalar tower \(R \to B_1 \to M\)), then \(M\) is module-finite over \(R\).
\end{lemma}
\begin{proof}

\leanok
  Compatibility of \(\psi^{-1}\) with the structure maps follows from that of \(\psi\), so
  \(\psi^{-1}\) is an \(R\)-algebra isomorphism; transporting \(R\)-finiteness of \(B_2\)
  across it makes \(B_1\) module-finite over \(R\). Transitivity of module-finiteness through
  the scalar tower \(R \to B_1 \to M\) then gives \(M\) module-finite over \(R\).
\end{proof}

\begin{lemma}\leanok
[Descent of a localised-module structure along a scalar tower]
  \label{lem:gf_islocalizedmodule_restrictscalars}
  \lean{AlgebraicGeometry.GenericFreeness.isLocalizedModule_restrictScalars}
  Let \(A_g \to R\) be a ring map with compatible \(A_g\)- and \(R\)-module structures on
  \(M\) and \(M'\) (scalar towers \(A_g \to R \to M\) and \(A_g \to R \to M'\)), and let
  \(S \subseteq A_g\) be a multiplicative subset. If an \(R\)-linear map
  \(f : M \to M'\) realises \(M'\) as the localisation of \(M\) at the image submonoid
  \(S \cdot 1 \subseteq R\) over \(R\), then its restriction of scalars to \(A_g\) realises
  \(M'\) as the localisation of \(M\) at \(S\) over \(A_g\).
\end{lemma}
\begin{proof}

\leanok
  One verifies the universal property of a localised module for the restricted map. Each
  \(s \in S\) acts invertibly on \(M'\): its action coincides, via the scalar tower, with
  that of its image in \(R\), which is a unit because \(f\) is a localisation at the image
  submonoid. Surjectivity onto \(M'\) by fractions \(x/s\) with \(s \in S\) and the
  equality criterion likewise descend from the corresponding image-submonoid statements,
  rewriting each \(R\)-scalar by an \(A_g\)-scalar through the scalar tower. Assembling
  these conditions gives the claim.
\end{proof}

\begin{lemma}[Localisation identification across an associated denominator]
  \label{lem:isLocalization_away_mul_of_associated}
  \lean{IsLocalization.Away.mul_of_associated}
  \mathlibok
  \textit{Provided by Mathlib.}
  Let \(R \to S \to T\) be a tower of commutative rings with compatible structure maps, let
  \(x, z \in R\) and \(y \in S\), and suppose \(S\) is the localisation of \(R\) away from
  \(x\) and \(T\) is the localisation of \(S\) away from \(y\). If the image of \(z\) in \(S\)
  is associated to \(y\) (the two differ by a unit), then \(T\) is the localisation of \(R\)
  away from the single product \(x z\); that is, \(T \cong R[(x z)^{-1}]\).
\end{lemma}

\begin{lemma}[Freeness transports across a semilinear ring isomorphism]
  \label{lem:module_free_of_ringEquiv}
  \lean{Module.Free.of_ringEquiv}
  \mathlibok
  \textit{Provided by Mathlib.}
  Let \(R, R'\) be rings, \(M\) an \(R\)-module and \(M'\) an \(R'\)-module, and let
  \(e : M \xrightarrow{\sim} M'\) be an additive isomorphism that is semilinear over a ring
  isomorphism \(\sigma : R \xrightarrow{\sim} R'\), i.e.\ \(e(r \cdot m) = \sigma(r) \cdot e(m)\).
  If \(M\) is free over \(R\), then \(M'\) is free over \(R'\).
\end{lemma}

\begin{lemma}\leanok
[L5 step-4 helper — descent of freeness across an iterated localisation]
  \label{lem:gf_away_tower_descent}
  \lean{AlgebraicGeometry.GenericFreeness.free_localizationAway_of_away_tower}
  \uses{lem:isLocalization_away_mul_of_associated, lem:module_free_of_ringEquiv}
  Let \(A\) be a domain and \(T\) an \(A\)-module. Let \(g \in A\), \(g \ne 0\), and let
  \(h \in A_g = A[g^{-1}]\) be non-zero. Write
  \((T_g)_h\) for the iterated localisation of \(T\) — first away from \(g\), then away from
  \(h\) — and \((A_g)_h = A_g[h^{-1}]\) for the corresponding iterated localisation of the
  base ring. If \((T_g)_h\) is free over \((A_g)_h\), then there exists \(f \in A\),
  \(f \ne 0\), such that the single localisation \(T_f\) is free over \(A_f = A[f^{-1}]\).
  Concretely the witness is the single product \(f := g\,a\) for an element \(a \in A\) whose
  image in \(A_g\) is associated to \(h\).
\end{lemma}

\begin{proof}
  \uses{lem:isLocalization_away_mul_of_associated, lem:module_free_of_ringEquiv}
  \leanok
  \emph{Clear the denominator of \(h\).} As \(h\) lies in \(A_g = A[g^{-1}]\),
  surjectivity of the localisation map \(A \to A_g\) onto fractions yields \(a \in A\) and a
  power \(s\) of \(g\) with \(h \cdot \bar s = \bar a\) in \(A_g\), where \(\bar{\,\cdot\,}\)
  denotes the image in \(A_g\). The image \(\bar s\) is a unit (a power of \(g\) is invertible
  in \(A_g\)), so \(\bar a = h\,\bar s\) is associated to \(h\); in particular \(\bar a \ne 0\)
  (because \(h \ne 0\)), hence \(a \ne 0\). Since \(A\) is a domain and \(g, a \ne 0\), the
  product \(f := g\,a\) is non-zero.

  \emph{Identify the base rings.} The associatedness of \(\bar a\) and \(h\), together with the
  tower \(A \to A_g \to A_h\), identifies the iterated localisation
  \((A_g)_h = A_g[h^{-1}]\) with the single localisation of \(A\) away from the product
  \(g\,a\): by \cref{lem:isLocalization_away_mul_of_associated} (with \(x = g\), \(z = a\),
  \(y = h\)) the ring \(A_h\) is a localisation of \(A\) away from \(g\,a\), and uniqueness of
  localisation furnishes a ring isomorphism \(A_{ga} \xrightarrow{\sim} A_h\).

  \emph{Identify the modules.} The composite localisation map
  \[
    \psi : T \longrightarrow (T_g)_h, \qquad
    \psi := \mathrm{mk}_{\langle h\rangle} \circ \mathrm{mk}_{\langle g\rangle},
  \]
  first into \(T_g\) and then into \((T_g)_h\), realises \((T_g)_h\) as the localisation of
  \(T\) at the powers of \(g\,a\): every \(g^n a^m\) acts invertibly on \((T_g)_h\), \(\psi\) is
  surjective up to such denominators, and its kernel is the expected one. This
  localisation-of-localisation identity is verified directly from the defining properties of a
  localised module; no packaged lemma supplies it.

  \emph{Transport freeness.} By hypothesis \((T_g)_h\) is free over \(A_h\). The ring
  isomorphism \(A_{ga} \cong A_h\) and the induced semilinear module equivalence
  \(T_{ga} \simeq (T_g)_h\) — assembled from the localised-module identification of the
  previous step and promoted to an \(A_{ga}\)-linear equivalence — carry this
  freeness back by \cref{lem:module_free_of_ringEquiv}: \(T_{ga}\) is free over
  \(A_{ga}\). Taking \(f := g\,a\) gives \(T_f\) free over \(A_f\), as required.
\end{proof}

\begin{lemma}\leanok
[L5 — polynomial-ring core of generic freeness]
  \label{lem:gf_polynomial_core}
  \lean{AlgebraicGeometry.GenericFreeness.exists_free_localizationAway_polynomial}
  \uses{lem:gf_finite_module, lem:gf_splice_shortExact, lem:gf_torsion_base,
        lem:gf_splice_shortExact_free_transport, lem:gf_generic_rank_ses,
        lem:gf_torsion_reindex, lem:gf_away_tower_descent}
  % SOURCE: [Nitsure], \S4 ``Lemma on Generic Flatness'', proof, inductive step
  % (read from references/nitsure-hilbert-quot-src/nitsure-hilbert-quot.tex, L1760-L1772).
  % SOURCE QUOTE PROOF:
  % "Let $m$ be the generic rank of the finite module $B_g$ over the domain
  %  $A_g[b_1,\ldots, b_n]$. Then we have a short exact sequence of
  %  $A_g[b_1,\ldots, b_n]$-modules of the form
  %  $0\to A_g[b_1,\ldots, b_n]^{\oplus m} \to B_g \to T \to 0$ where $T$ is a finite
  %  torsion module over $A_g[b_1,\ldots, b_n]$. Therefore, the dimension of the
  %  support of $K\otimes_{A_g}T$ as a $K\otimes_{A_g}(B_g)$-module is strictly less
  %  than $n$. Hence by induction on $n$ (applied to the data $A_g$, $B_g$, $T$),
  %  there exists some $h\ne 0$ in $A$ with $T_h$ free over $A_{gh}$. Taking $f=gh$,
  %  the lemma follows from the above short exact sequence."
  \textit{Source: [Nitsure], \S4, ``Lemma on Generic Flatness'' (polynomial-ring
  core of the induction).}
  Let \(A\) be a noetherian domain and \(d \ge 0\). Let \(N\) be a finite module over
  the polynomial ring \(A[X_1, \dots, X_d]\), regarded as an \(A\)-module. Then there
  exists \(f \in A\), \(f \ne 0\), such that \(N_f\) is free over \(A_f\).

\end{lemma}

\begin{proof}
  \uses{lem:gf_finite_module, lem:gf_splice_shortExact, lem:gf_torsion_base,
        lem:gf_splice_shortExact_free_transport, lem:gf_generic_rank_ses,
        lem:gf_torsion_reindex, lem:gf_away_tower_descent}
        \leanok
  % SOURCE QUOTE PROOF: see the comment above the statement (L1760-L1772).
  \textit{Source: [Nitsure], \S4 (polynomial-ring core).}
  We argue by strong induction on \(d\), \emph{generalising the base domain \(A\)} (along
  with the module \(N\)): the inductive hypothesis is available for every \(d' < d\),
  every noetherian domain in place of \(A\), and every finite module over the
  corresponding polynomial ring.

  \emph{Base case \(d = 0\).} The ring \(P_0 = A[X_1, \dots, X_0]\) is \(A\) itself, so
  \(N\) is a finite \(A\)-module and the claim is exactly the finite-module leaf
  \cref{lem:gf_finite_module}.

  \emph{Inductive step \(d \ge 1\).} If \(N\) is already a torsion \(P_d\)-module over
  \(A\) — equivalently \(N_K = 0\) for \(K = \mathrm{Frac}(A)\) — then the torsion base
  case \cref{lem:gf_torsion_base} produces \(f \ne 0\) with \(N_f = 0\), which is free
  over \(A_f\). Otherwise we proceed in five steps.
  \begin{enumerate}
    \item \textbf{Generic-rank SES.} Apply \cref{lem:gf_generic_rank_ses} to \(N\) over
      \(P_d = A[X_1, \dots, X_d]\): there are \(m \in \mathbb{N}\) and an injective
      \(P_d\)-linear map \(\varphi : P_d^{\oplus m} \to N\) with torsion cokernel
      \(T := N / \operatorname{im}\varphi\), i.e.\ a short exact sequence of
      \(P_d\)-modules
      \[
        0 \to P_d^{\oplus m} \xrightarrow{\ \varphi\ } N \to T \to 0,
        \qquad T \text{ torsion}.
      \]
      No element of \(A\) is inverted at this step.
    \item \textbf{Reindex the torsion cokernel.} Apply \cref{lem:gf_torsion_reindex} to the
      finite torsion module \(T\): there are \(g \ne 0\) in \(A\) and \(m' < d\) such that
      the localisation \(T_g\) is module-finite over the polynomial ring
      \(A_g[X_1, \dots, X_{m'}]\) in fewer variables over the base ring \(A_g = A[g^{-1}]\).
    \item \textbf{Inductive hypothesis at the new base \(A_g\).} Since \(A_g\) is again a
      noetherian domain and \(m' < d\), the inductive hypothesis applies \emph{at base
      \(A_g\)} to the finite \(A_g[X_1, \dots, X_{m'}]\)-module \(T_g\): there is
      \(h \ne 0\) in \(A_g\) with \((T_g)_h\) free over \((A_g)_h\). This is exactly the
      point at which the base ring of the induction must be allowed to change from \(A\)
      to \(A_g\).
    \item \textbf{Descend the witness from \(A_g\) to \(A\).} Clearing the denominator of
      \(h\) over \(A_g\) produces a numerator \(a \in A\) whose image in \(A_g\) is associated
      to \(h\); then the iterated localisation \((A_g)_h\) is the single localisation
      \(A_{g a}\) and \((T_g)_h\) is \(T_{g a}\). Thus freeness of \((T_g)_h\) over
      \((A_g)_h\) transports to freeness of \(T_{g a}\) over \(A_{g a}\) by the
      iterated-\(\mathrm{Away}\)-tower descent \cref{lem:gf_away_tower_descent}. Set
      \(f := g a \ne 0\) — the single product, \emph{not} a square — so that \(T_f\) is free
      over \(A_f\).
    \item \textbf{Splice.} The left end \(P_d^{\oplus m}\) of the sequence is a free
      \(A\)-module (the polynomial ring \(P_d = A[X_1, \dots, X_d]\) is \(A\)-free, hence
      so is any finite direct power; note \(P_d^{\oplus m}\) is \emph{not} module-finite
      over \(A\), but freeness is all the splice requires), and a free module stays free
      under the localisation \(A \to A_f\); so \(P_d^{\oplus m}\) localises to a free
      \(A_f\)-module. With both ends free over \(A_f\) after inverting \(f = g a\),
      \cref{lem:gf_splice_shortExact} splices the localised short exact sequence to
      conclude that the middle term \(N_f\) is free over \(A_f\). Taking \(f = g a\)
      completes the induction.
  \end{enumerate}
\end{proof}

\subsection{Supporting finite-module helpers}
\label{sec:gf_finite_helpers}

The following two lemmas package the finite \(A\)-module leaf in the two forms
(flatness and module-finiteness) used by the dévissage chain.

\begin{lemma}



\leanok
[Flat localization away, finite module]
  \label{lem:gf_flat_finite}
  \lean{AlgebraicGeometry.GenericFreeness.exists_flat_localizationAway_of_finite}
  \uses{lem:gf_finite_module}
  Let \(A\) be a noetherian domain and \(M\) a finite \(A\)-module. Then there exists
  \(f \in A\), \(f \ne 0\), such that \(M_f\) is flat over \(A_f\).
\end{lemma}
\begin{proof}
  \uses{lem:gf_finite_module}
  \leanok
  By \cref{lem:gf_finite_module} \(M_f\) is free over \(A_f\); free modules are flat.
\end{proof}

\begin{lemma}



\leanok
[Free localization away, module-finite variant]
  \label{lem:gf_free_moduleFinite}
  \lean{AlgebraicGeometry.GenericFreeness.exists_free_localizationAway_of_moduleFinite}
  \uses{lem:gf_finite_module}
  Let \(A\) be a noetherian domain, \(B\) a module-finite \(A\)-algebra, and \(M\) a
  finite \(B\)-module with the compatible \(A\)-module structure via the scalar tower
  \(A \to B \to M\). Then there exists \(f \in A\), \(f \ne 0\), such that \(M_f\)
  is free over \(A_f\).
\end{lemma}
\begin{proof}
  \uses{lem:gf_finite_module}
  \leanok
  Since \(B\) is module-finite over \(A\) and \(M\) is module-finite over \(B\),
  transitivity of module-finiteness makes \(M\) a finite \(A\)-module.
  \cref{lem:gf_finite_module} then gives the required \(f\).
\end{proof}

\subsection{Geometric bridges for the affine assembly}
\label{sec:gf_geometric_bridges}

The geometric form of generic flatness is assembled from the algebraic core by
restricting to a non-empty affine open of \(S\) and working patchwise over a finite
affine cover of the preimage. Two genuinely geometric steps are isolated here as
named bridges: the affine-local finiteness of the section module of a finite-type
quasi-coherent sheaf (G1), feeding the algebraic core its finite-module input; and
the flat-locality assembly (G3), promoting per-patch freeness to flatness of the
restricted sheaf over the base open.

\begin{lemma}[Finiteness is local on a finite spanning family]
  \label{lem:module_finite_of_localizationSpan}
  \lean{Module.Finite.of_localizationSpan_finite}
  \mathlibok
  \textit{Provided by Mathlib (\texttt{Mathlib.RingTheory.LocalProperties.Submodule}).}
  Let \(R\) be a commutative ring, \(N\) an \(R\)-module, and \(t \subseteq R\) a
  finite family with \(\mathrm{span}\,t = \top\) (the \(t\) generate the unit
  ideal). Suppose that for each \(g \in t\) the localization \(N_g\) is a finite
  module over the localization \(R_g\). Then \(N\) is a finite \(R\)-module.
\end{lemma}

\begin{lemma}\leanok
[Half 1 helper — transfer of localized-module finiteness across models]
  \label{lem:finite_localizedModule_transfer}
  \lean{AlgebraicGeometry.finite_localizedModule_of_isLocalizedModule}
  Let \(R\) be a commutative ring, \(S \subseteq R\) a multiplicative submonoid,
  and \(M\) an \(R\)-module. Let \(R_S\) be an \(R\)-algebra that localizes \(R\)
  at \(S\), and let \(N\) be an \(R\)-module carrying a compatible \(R_S\)-module
  structure, together with an \(R\)-linear map \(\varphi : M \to N\) exhibiting
  \(N\) as the localization of \(M\) at \(S\) (\(\mathrm{IsLocalizedModule}\,S\,
  \varphi\)). If \(N\) is a finite \(R_S\)-module, then the canonical localized
  module \(\mathrm{LocalizedModule}\,S\,M\) is a finite module over the canonical
  localization \(\mathrm{Localization}\,S\).
\end{lemma}

\begin{proof}

\leanok
  The localization map \(\varphi\) provides an \(R\)-linear isomorphism
  \(e : N \xrightarrow{\sim} \mathrm{LocalizedModule}\,S\,M\) (universal property
  of localized modules), and the algebra localization provides an \(R\)-algebra
  isomorphism \(\psi : R_S \xrightarrow{\sim} \mathrm{Localization}\,S\). These
  are compatible: \(e\) is semilinear over \(\psi\), i.e.\ \(e(a \cdot x) =
  \psi(a) \cdot e(x)\) for \(a \in R_S\), \(x \in N\) (both sides are pinned by
  clearing a common denominator \(s \in S\), where the \(R\)-actions agree). A
  finite \(R_S\)-spanning set \(T\) of \(N\) therefore maps under \(e\) to a
  finite set whose \(\mathrm{Localization}\,S\)-span is all of
  \(\mathrm{LocalizedModule}\,S\,M\); hence the canonical model is finite over
  \(\mathrm{Localization}\,S\).
\end{proof}

\begin{lemma}\leanok
[Half 1 — G1 locality reduction over a finite basic-open cover]
  \label{lem:gf_finite_sections_of_basicOpen_finite_cover}
  \lean{AlgebraicGeometry.gf_finite_sections_of_basicOpen_finite_cover}
  \uses{lem:qcoh_section_localization_basicOpen, lem:finite_localizedModule_transfer,
        lem:module_finite_of_localizationSpan, lem:isLocalization_basicOpen_mathlib}
  % SOURCE: [Stacks Project], "Properties of Schemes", Tag 01PB, proof, gluing step
  % (read from references/stacks-properties.tex, L2107-L2109).
  % SOURCE QUOTE PROOF:
  % "Thus $M_{f_i}$ is finitely generated for each $i$. Hence we conclude by
  %  Algebra, Lemma \ref{algebra-lemma-cover}."
  \textit{Source: [Stacks Project], "Properties of Schemes", Tag 01PB (gluing step,
  Algebra Lemma \texttt{lemma-cover}).}
  Let \(X\) be a scheme, \(\F\) a quasi-coherent sheaf of modules on \(X\), and
  \(W \subseteq X\) an affine open. Suppose \(t \subseteq \Gamma(X, W)\) is a finite
  family with \(\mathrm{Ideal.span}\,t = \top\) (so the basic opens \(D(g)\),
  \(g \in t\), cover \(W\)), and that for each \(g \in t\) the section module
  \(\Gamma(\F, D(g))\) is finite over \(\Gamma(X, D(g))\). Then \(\Gamma(\F, W)\)
  is a finite \(\Gamma(X, W)\)-module.
\end{lemma}

\begin{proof}
  \uses{lem:qcoh_section_localization_basicOpen, lem:finite_localizedModule_transfer,
        lem:module_finite_of_localizationSpan, lem:isLocalization_basicOpen_mathlib}
        \leanok
  \textit{Source: [Stacks Project], Tag 01PB (gluing step).}
  By the local-on-a-spanning-family criterion \cref{lem:module_finite_of_localizationSpan}
  it suffices to show, for each \(g \in t\), that the localization of
  \(\Gamma(\F, W)\) at \(g\) is finite over the localization
  \((\Gamma(X, W))_g\). Fix \(g \in t\). The basic open \(D(g)\) is affine, and the
  restriction \(\Gamma(X, W) \to \Gamma(X, D(g))\) exhibits \(\Gamma(X, D(g))\) as
  the away-localization \((\Gamma(X, W))_g\)
  (\cref{lem:isLocalization_basicOpen_mathlib}); moreover the section restriction
  \(\Gamma(\F, W) \to \Gamma(\F, D(g))\) exhibits \(\Gamma(\F, D(g))\) as the
  localized module \((\mathrm{powers}\,g)^{-1}\,\Gamma(\F, W)\)
  (\cref{lem:qcoh_section_localization_basicOpen}). Thus \(\Gamma(\F, D(g))\) is a
  model of the localized module, finite over \(\Gamma(X, D(g)) = (\Gamma(X, W))_g\)
  by hypothesis; the model-transfer lemma
  \cref{lem:finite_localizedModule_transfer} carries this finiteness to the
  canonical localized module of \(\Gamma(\F, W)\) at the powers of \(g\), over
  \((\Gamma(X, W))_g\). This supplies the per-element hypothesis, and
  \cref{lem:module_finite_of_localizationSpan} concludes that \(\Gamma(\F, W)\) is
  finite over \(\Gamma(X, W)\).
\end{proof}

\subsection{Transport engine for generating sections}

The restriction-of-generators lemma below is assembled from a generic transport
principle: a finite generating family of a sheaf of modules \(\mathcal{M}\) pushes
forward to a finite generating family of \(F(\mathcal{M})\) along any functor \(F\) of
sheaves of modules that preserves colimits and carries the structure sheaf to the
structure sheaf. The three blocks of this subsection package that principle (the
generating family, its index type, and the finiteness bookkeeping). They are
project-bespoke: Mathlib provides the analogous transport only for full presentations
(with kernel data), not for bare generating families.

\begin{lemma}\leanok
[Transport of a generating family along a colimit-preserving functor]
  \label{lem:genSections_map}
  \lean{AlgebraicGeometry.SheafOfModules.GeneratingSections.map}
  Let \(F : \mathrm{Mod}\,R \to \mathrm{Mod}\,S\) be a functor of sheaves of modules
  that preserves colimits and is equipped with an isomorphism
  \(\eta : F(\OO_R) \xrightarrow{\sim} \OO_S\) carrying the unit (structure sheaf) of the
  source to that of the target. Given a generating family \(\sigma\) of \(\mathcal{M}\)
  with index set \(I = \sigma.I\) and epimorphism
  \(\sigma.\pi : \OO_R^{\oplus I} \twoheadrightarrow \mathcal{M}\), the images
  \((\mathrm{mapFree}\,F\,\eta\,I)^{-1} \circ F(\sigma.\pi)\) form a generating family of
  \(F(\mathcal{M})\) with the same index set \(I\). The map is an epimorphism because it
  is the composite of an isomorphism with \(F(\sigma.\pi)\), and \(F\) preserves the
  epimorphism \(\sigma.\pi\).
\end{lemma}

\begin{lemma}\leanok
[The transported family keeps its index type]
  \label{lem:genSections_map_I}
  \lean{AlgebraicGeometry.SheafOfModules.GeneratingSections.map_I}
  \uses{lem:genSections_map}
  The index set of the transported generating family of \cref{lem:genSections_map} equals
  the index set of the original: \((\mathrm{map}\,\sigma\,F\,\eta)\!.I = \sigma.I\). This
  holds definitionally.
\end{lemma}

\begin{lemma}\leanok
[Finiteness survives transport]
  \label{lem:genSections_map_isFiniteType}
  \lean{AlgebraicGeometry.SheafOfModules.GeneratingSections.map_isFiniteType}
  \uses{lem:genSections_map, lem:genSections_map_I}
  If the original generating family \(\sigma\) is of finite type (its index set is finite),
  then so is the transported family \(\mathrm{map}\,\sigma\,F\,\eta\): the index type is
  unchanged by \cref{lem:genSections_map_I}, hence still finite.
\end{lemma}

\begin{lemma}\leanok
[Half 2, seam 1a — local generators restrict to a smaller open]
  \label{lem:gf_localGenerators_restrict}
  \lean{AlgebraicGeometry.gf_localGenerators_restrict}
  \uses{lem:genSections_map, lem:genSections_map_I, lem:genSections_map_isFiniteType}
  Let \(X\) be a scheme and \(\F\) a quasi-coherent sheaf of \(\mathcal{O}_X\)-modules.
  Suppose that on an open \(Y \subseteq X\) the restriction \(\F|_Y\) is
  globally generated by finitely many sections. Then for any open \(V \le Y\) the
  restricted family — the images of those generators under the section-restriction
  \(\Gamma(\F, Y) \to \Gamma(\F, V)\) — is again a finite generating family of
  \(\F|_V\). Generation survives passage to the smaller open (a family that generates
  every stalk over \(Y\) generates every stalk over \(V\)), and the index set, and hence
  finiteness, is unchanged.
\end{lemma}

\begin{proof}
  \uses{lem:genSections_map, lem:genSections_map_I, lem:genSections_map_isFiniteType}
  \leanok
  A family of sections \(\{s_i\}_{i \in I}\) generates \(\F|_Y\) precisely when the
  induced map \(\OO_Y^{\oplus I} \to \F|_Y\) is surjective on every stalk
  \(\F_y\,(y \in Y)\). Restricting along the inclusion \(V \le Y\) the
  generators map to \(\{s_i|_V\}_{i \in I}\), and for a point \(x \in V\) the stalk
  \(\F_x\) and the comparison map agree with the value at \(x\) computed over \(Y\), since
  stalks are insensitive to shrinking the ambient open. Hence the induced map
  \(\OO_V^{\oplus I} \to \F|_V\) is surjective on every stalk over \(V\), so
  \(\{s_i|_V\}_{i \in I}\) generates \(\F|_V\). The index set \(I\) is the same
  finite set, so finiteness is preserved. Concretely the restricted family is obtained by
  transporting the generating family along the (colimit-preserving, unit-preserving)
  restriction functors via the transport engine \cref{lem:genSections_map}, whose index
  type is unchanged (\cref{lem:genSections_map_I}) and which preserves finiteness
  (\cref{lem:genSections_map_isFiniteType}).
\end{proof}

\begin{lemma}\leanok
[Half 2, seam 1b — finite standard-basic-open subcover of an affine open]
  \label{lem:gf_affine_finite_standard_subcover}
  \lean{AlgebraicGeometry.gf_affine_finite_standard_subcover}
  Let \(W\) be an affine open of a scheme \(X\), with the canonical identification
  \(W \cong \Spec \Gamma(X, W)\). Given an arbitrary open cover \(\{U_\alpha\}\) of \(W\),
  there is a finite family \(t \subseteq \Gamma(X, W)\) with
  \(\mathrm{Ideal.span}\,t = \top\) such that each basic open \(D(g)\) (\(g \in t\)) is
  contained in some cover member \(U_\alpha\). In particular the \(D(g)\,(g \in t)\) form a
  finite standard-basic-open subcover of \(W\) refining \(\{U_\alpha\}\).
\end{lemma}

\begin{proof}

\leanok
  Under \(W \cong \Spec \Gamma(X, W)\) the basic opens \(D(g)\,(g \in \Gamma(X, W))\) form a
  basis of the topology (\(\mathrm{PrimeSpectrum.isBasis\_basic\_opens}\)). Hence each
  cover member \(U_\alpha\) is a union of basic opens, and the collection of all basic
  opens \(D(g)\) that are contained in some \(U_\alpha\) still covers \(W\). The affine
  scheme \(W\) is quasi-compact, so finitely many of these basic opens
  \(D(g_1), \dots, D(g_n)\) already cover \(W\); set
  \(t = \{g_1, \dots, g_n\} \subseteq \Gamma(X, W)\), a finite family, each \(D(g_i)\)
  contained in some \(U_\alpha\) by construction. That the \(D(g_i)\) cover
  \(W = \Spec \Gamma(X, W)\) — i.e. no prime ideal contains all the \(g_i\) — is exactly
  the condition \(\mathrm{Ideal.span}\,t = \top\).
\end{proof}

\begin{lemma}[Half 2, seam 1c — finite generation equals a finite free epimorphism]
  \label{lem:gf_finite_gen_iff_free_epi}
  \lean{AlgebraicGeometry.gf_finite_gen_iff_free_epi}
  Let \(Y\) be an open of a scheme and \(\F\) a sheaf of modules on \(Y\). Then
  \(\F\) is globally generated by a finite family of sections if and only if there is a
  finite index set \(I\) and an epimorphism of \(\OO_Y\)-modules
  \(\OO_Y^{\oplus I} \twoheadrightarrow \F\). No quasi-coherence of \(\F\) is assumed:
  the statement holds for arbitrary sheaves of modules on \(Y\), and applies in particular
  to the restrictions \(\F|_{D(g)}\) arising in the cover assembly.
\end{lemma}

\begin{proof}
  A finite family of sections \(\{s_i\}_{i \in I}\) of \(\F|_Y\) is the same datum as an
  \(\OO_Y\)-linear map \(\OO_Y^{\oplus I} \to \F|_Y\) sending the \(i\)-th basis section to
  \(s_i\). The family generates \(\F|_Y\) (the induced map is surjective on every stalk)
  exactly when this map is an epimorphism of sheaves of modules. Thus a finite generating
  family yields a finite-free epimorphism, and conversely the images of the basis sections
  of any epimorphism \(\OO_Y^{\oplus I} \twoheadrightarrow \F|_Y\) with \(I\) finite form a
  finite generating family.
\end{proof}

\begin{lemma}\leanok
[Half 2, seam 1 — finite type yields a finite cover of globally-generated affines]
  \label{lem:gf_finiteType_affine_finite_cover_generated}
  \lean{AlgebraicGeometry.gf_finiteType_affine_finite_cover_generated}
  % SOURCE: [Stacks Project], "Properties of Schemes", Tag 01PB, proof
  % (read from references/stacks-properties.tex, L2101-L2107).
  % SOURCE QUOTE PROOF:
  % "Assume $\widetilde M$ is a finite type $\mathcal{O}_X$-module. This means
  %  there exists an open covering of $X$ such that $\widetilde M$ restricted to
  %  the members of this covering is globally generated by finitely many sections.
  %  Thus there also exists a standard open covering
  %  $X = \bigcup_{i = 1, \ldots, n} D(f_i)$ such that $\widetilde M|_{D(f_i)}$
  %  is generated by finitely many sections."
  \uses{lem:gf_localGenerators_restrict, lem:gf_affine_finite_standard_subcover,
        lem:gf_finite_gen_iff_free_epi}
  \textit{Source: [Stacks Project], "Properties of Schemes", Tag 01PB (standard
  cover of generated affines).}
  Let \(X\) be a scheme, \(\F\) a sheaf of modules of finite type,
  and \(W \subseteq X\) an affine open. Then there is a finite family
  \(t \subseteq \Gamma(X, W)\) with \(\mathrm{Ideal.span}\,t = \top\) such that for
  each \(g \in t\) the basic open \(D(g)\) is affine and \(\F|_{D(g)}\) is globally
  generated by finitely many sections; equivalently, for each \(g \in t\) there is
  a finite index set \(I_g\) and an epimorphism of \(\OO_{D(g)}\)-modules
  \(\OO_{D(g)}^{\oplus I_g} \twoheadrightarrow \F|_{D(g)}\).
\end{lemma}

\begin{proof}
  \uses{lem:gf_localGenerators_restrict, lem:gf_affine_finite_standard_subcover,
        lem:gf_finite_gen_iff_free_epi}
        \leanok
  \textit{Source: [Stacks Project], Tag 01PB (standard cover of generated affines).}
  Finite type of \(\F\) means that \(X\) has an open covering \(\{Y_\alpha\}\) on each
  member of which \(\F\) carries a finite generating family of sections. The intersections
  \(\{Y_\alpha \cap W\}\) form an open cover of the affine \(W\). By the finite standard
  subcover refinement \cref{lem:gf_affine_finite_standard_subcover} there is a finite family
  \(t = \{g_1, \dots, g_n\} \subseteq \Gamma(X, W)\) with
  \(\mathrm{Ideal.span}\,t = \top\) and each basic open \(D(g_i)\) contained in some
  \(Y_\alpha \cap W\) — in particular \(D(g_i) \subseteq W\) is affine, and \(D(g_i)\)
  embeds by an open immersion into the member \(Y_\alpha\) on which \(\F\) is finitely
  generated. By the restriction-of-generators lemma
  \cref{lem:gf_localGenerators_restrict} the finite generating family on \(Y_\alpha\)
  restricts to a finite generating family of \(\F|_{D(g_i)}\), so \(\F|_{D(g_i)}\) is
  globally generated by finitely many sections. Finally, by the free-epimorphism
  reformulation \cref{lem:gf_finite_gen_iff_free_epi}, this is equivalent to the existence
  of a finite index set \(I_{g_i}\) and an epimorphism of \(\OO_{D(g_i)}\)-modules
  \(\OO_{D(g_i)}^{\oplus I_{g_i}} \twoheadrightarrow \F|_{D(g_i)}\). That the \(D(g_i)\)
  cover \(W = \Spec \Gamma(X, W)\) is exactly \(\mathrm{Ideal.span}\,t = \top\). This is
  Stacks' passage from the local-generation datum to a standard finite cover of generated
  affines.
\end{proof}

\begin{lemma}[Affine $\widetilde{(-)}\dashv\Gamma$ adjunction with its $\mathrm{fromTildeΓ}$ counit]
  \label{lem:tilde_adjunction_mathlib}
  \lean{AlgebraicGeometry.tilde.adjunction}
  \mathlibok
  \textit{Provided by Mathlib (\texttt{Mathlib.AlgebraicGeometry.Modules.Tilde}).}
  On the affine scheme \(\Spec R\) the associated-sheaf functor
  \(\widetilde{(-)} : \mathrm{Mod}\,R \to (\Spec R).\mathrm{Modules}\) is left adjoint to the
  affine section functor \(\Gamma : (\Spec R).\mathrm{Modules} \to \mathrm{Mod}\,R\) (Mathlib's
  \(\mathrm{moduleSpecΓFunctor}\), evaluation of a sheaf of modules at the top open). The counit of
  this adjunction is the natural transformation
  \(\Gamma(-)\,\widetilde{\phantom{x}} \to \mathrm{id}\) whose component at a sheaf \(\mathcal{M}\)
  is \(\mathcal{M}.\mathrm{fromTildeΓ} : \widetilde{\Gamma(\mathcal{M}, \top)} \to \mathcal{M}\), and
  the left adjoint \(\widetilde{(-)}\) is faithful.
\end{lemma}

\begin{lemma}\leanok
[Half 2, seam 2 — $\Gamma$ on affine quasi-coherent modules sends epis to surjections]
  \label{lem:gf_affine_qcoh_Gamma_epi}
  \lean{AlgebraicGeometry.gf_affine_qcoh_Gamma_epi}
  \uses{lem:qcoh_affine_section_localization, lem:tilde_adjunction_mathlib,
    lem:isIso_fromTildeΓ_of_isLocalizedModule_restrict}
  % SOURCE: [Stacks Project], "Properties of Schemes", Tag 01PB, proof
  % (read from references/stacks-properties.tex, L2106-L2107) — the implication
  % "$\widetilde M|_{D(f_i)}$ generated by finitely many sections $\Rightarrow$
  % $M_{f_i}$ finitely generated" is exactly section-level surjectivity of a sheaf
  % epi over an affine.
  % SOURCE QUOTE PROOF:
  % "such that $\widetilde M|_{D(f_i)}$ is generated by finitely many sections.
  %  Thus $M_{f_i}$ is finitely generated for each $i$."
  \textit{Source: [Stacks Project], "Properties of Schemes", Tag 01PB (sections of a
  generated sheaf over an affine).}
  Let \(V\) be an affine open of a scheme, with \(B = \Gamma(X, V)\) and
  \(V \cong \Spec B\), and let \(\pi : \mathcal{G} \to \F\) be an epimorphism of
  quasi-coherent sheaves of modules over \(V\). Then the induced map on sections
  \(\Gamma(\pi) : \Gamma(\mathcal{G}, V) \to \Gamma(\F, V)\) is a surjective
  \(B\)-module map.
\end{lemma}

\begin{proof}
  \uses{lem:qcoh_affine_section_localization, lem:tilde_adjunction_mathlib,
    lem:isIso_fromTildeΓ_of_isLocalizedModule_restrict}
    \leanok
  \textit{Source: [Stacks Project], Tag 01PB (sections of a generated sheaf over an
  affine).}
  Write \(V \cong \Spec B\) and let \(\Gamma\) be the affine section functor on
  \((\Spec B).\mathrm{Modules}\), so that \(\Gamma(\pi) = \Gamma(\mathcal{G}, V) \to
  \Gamma(\F, V)\) is the map in question. By \cref{lem:tilde_adjunction_mathlib} the
  associated-sheaf functor \(\widetilde{(-)}\) is left adjoint to \(\Gamma\), with counit the
  natural transformation whose component at a sheaf \(\mathcal{M}\) is
  \(\mathcal{M}.\mathrm{fromTildeΓ} : \widetilde{\Gamma(\mathcal{M}, V)} \to \mathcal{M}\).

  Because \(\mathcal{G}\) and \(\F\) are quasi-coherent, the qcoh-sections descent
  (\cref{lem:qcoh_affine_section_localization}, fed through the counit-iso assembly
  \cref{lem:isIso_fromTildeΓ_of_isLocalizedModule_restrict}) makes both counit components
  \(\mathcal{G}.\mathrm{fromTildeΓ}\) and \(\F.\mathrm{fromTildeΓ}\) isomorphisms. Naturality of the
  counit on \(\pi\) gives the commuting square
  \[
    \widetilde{\Gamma(\pi)} \circ \mathcal{G}.\mathrm{fromTildeΓ}^{-1}
      = \F.\mathrm{fromTildeΓ}^{-1} \circ \pi,
    \qquad\text{equivalently}\qquad
    \widetilde{\Gamma(\pi)}
      = \F.\mathrm{fromTildeΓ}^{-1} \circ \pi \circ \mathcal{G}.\mathrm{fromTildeΓ}.
  \]
  The right-hand side is a composite of an isomorphism, the epimorphism \(\pi\), and an
  isomorphism, hence is itself an epimorphism. Thus \(\widetilde{\Gamma(\pi)}\) is an epimorphism
  of sheaves of modules. Since the left adjoint \(\widetilde{(-)}\) is faithful, it reflects
  epimorphisms; therefore \(\Gamma(\pi)\) is an epimorphism in \(\mathrm{Mod}\,B\), i.e. the
  \(B\)-linear map \(\Gamma(\pi)\) is surjective.

  This is the structural content of the vanishing of higher quasi-coherent cohomology on an
  affine — the failure of \(\Gamma\) to be right exact would obstruct exactly this surjectivity —
  obtained here directly from the affine adjunction without invoking a cohomology computation.
\end{proof}

\begin{lemma}\leanok
[Half 2, seam 2$+$3 — a qcoh epimorphism with finite source sections has finite target sections]
  \label{lem:gf_qcoh_finite_sections_globally_generated}
  \lean{AlgebraicGeometry.gf_qcoh_finite_sections_globally_generated}
  \uses{lem:gf_affine_qcoh_Gamma_epi, lem:module_finite_of_surjective_mathlib}
  % SOURCE: [Stacks Project], "Properties of Schemes", Tag 01PB
  % (read from references/stacks-properties.tex, L2094-L2107).
  % SOURCE QUOTE:
  % "Let $X = \Spec(R)$ be an affine scheme. The quasi-coherent sheaf of
  %  $\mathcal{O}_X$-modules $\widetilde M$ is a finite type $\mathcal{O}_X$-module
  %  if and only if $M$ is a finite $R$-module."
  \textit{Source: [Stacks Project], "Properties of Schemes", Tag 01PB.}
  Let \(V\) be an affine open of a scheme \(X\), write \(B = \Gamma(X, V)\) with
  \(V \cong \Spec B\), and let \(\pi : \mathcal{G} \to \F\) be an epimorphism of
  quasi-coherent sheaves of modules over \(V\). Suppose the source has finite sections:
  \(\Gamma(\mathcal{G}, V)\) is a finite \(B\)-module. Then \(\Gamma(\F, V)\) is a finite
  \(B\)-module.
\end{lemma}

\begin{proof}
  \uses{lem:gf_affine_qcoh_Gamma_epi, lem:module_finite_of_surjective_mathlib}
  \leanok
  \textit{Source: [Stacks Project], Tag 01PB.}
  Apply seam 2 (\cref{lem:gf_affine_qcoh_Gamma_epi}) to the epimorphism
  \(\pi : \mathcal{G} \to \F\) over the affine \(V\): the section functor yields a
  surjective \(B\)-module map \(\Gamma(\pi) : \Gamma(\mathcal{G}, V) \to \Gamma(\F, V)\).
  The source \(\Gamma(\mathcal{G}, V)\) is a finite \(B\)-module by hypothesis. A finite
  module surjecting onto \(\Gamma(\F, V)\) forces \(\Gamma(\F, V)\) to be finite, by
  \cref{lem:module_finite_of_surjective_mathlib}. Hence \(\Gamma(\F, V)\) is a finite
  \(B\)-module.
\end{proof}

\begin{lemma}\leanok
[Half 2 base case — a finite-free qcoh epimorphism has finite affine sections]
  \label{lem:gf_qcoh_sections_free_epi}
  \lean{AlgebraicGeometry.gf_qcoh_finite_sections_of_free_epi}
  \uses{lem:gf_qcoh_finite_sections_globally_generated, lem:tilde_adjunction_mathlib}
  % SOURCE: [Stacks Project], "Properties of Schemes", Tag 01PB, proof
  % (read from references/stacks-properties.tex, L2106-L2107) — "$\widetilde M|_{D(f_i)}$
  % generated by finitely many sections $\Rightarrow$ $M_{f_i}$ finitely generated".
  % SOURCE QUOTE PROOF:
  % "such that $\widetilde M|_{D(f_i)}$ is generated by finitely many sections.
  %  Thus $M_{f_i}$ is finitely generated for each $i$."
  \textit{Source: [Stacks Project], "Properties of Schemes", Tag 01PB (free generation
  over an affine).}
  Let \(V\) be an affine open of a scheme \(X\), write \(R = \Gamma(X, V)\) with
  \(V \cong \Spec R\), and let \(\F\) be a quasi-coherent sheaf of modules. Suppose there
  is a finite \(R\)-module \(N\) and an epimorphism of \(\OO_V\)-modules
  \(\widetilde{N} \twoheadrightarrow \F\) out of the associated sheaf of \(N\) (the case
  \(N = R^{\oplus I}\), \(I\) finite, being the free-generation datum
  \(\OO_V^{\oplus I} \twoheadrightarrow \F\)). Then \(\Gamma(\F, V)\) is a finite
  \(R\)-module.
\end{lemma}

\begin{proof}
  \uses{lem:gf_qcoh_finite_sections_globally_generated, lem:tilde_adjunction_mathlib}
  \leanok
  \textit{Source: [Stacks Project], Tag 01PB.}
  This is the free-source specialisation of
  \cref{lem:gf_qcoh_finite_sections_globally_generated}; what remains is to discharge that
  lemma's two source-side hypotheses for \(\mathcal{G} = \widetilde{N}\). By the affine
  \(\widetilde{(-)}\dashv\Gamma\) adjunction \cref{lem:tilde_adjunction_mathlib} the unit
  at \(N\) is an isomorphism \(N \xrightarrow{\sim} \Gamma(\widetilde{N}, V)\) (the left
  adjoint \(\widetilde{(-)}\) is fully faithful), so \(\Gamma(\widetilde{N}, V)\) is a
  finite \(R\)-module because \(N\) is; and the counit component at \(\widetilde{N}\),
  namely \(\widetilde{N}.\mathrm{fromTildeΓ}\), is an isomorphism (a fully-faithful left
  adjoint has invertible counit on its image), supplying the source quasi-coherence
  needed by seam 2. With both source hypotheses in hand,
  \cref{lem:gf_qcoh_finite_sections_globally_generated} applied to
  \(\widetilde{N} \twoheadrightarrow \F\) yields that \(\Gamma(\F, V)\) is a finite
  \(R\)-module.
\end{proof}

\begin{lemma}

\leanok
[Finiteness transports across a ring isomorphism and a semilinear additive map]
  \label{lem:module_finite_of_ringEquiv_semilinear}
  \lean{AlgebraicGeometry.module_finite_of_ringEquiv_semilinear}
  Let \(\sigma : R \xrightarrow{\sim} R'\) be a ring isomorphism and
  \(e : M \xrightarrow{\sim} M'\) an additive isomorphism that is \(\sigma\)-semilinear,
  i.e. \(e(r \cdot m) = \sigma(r) \cdot e(m)\) for all \(r \in R\), \(m \in M\). If
  \(M\) is a finite \(R\)-module then \(M'\) is a finite \(R'\)-module.
\end{lemma}

\begin{proof}

\leanok
  Pick a finite \(R\)-spanning set \(T \subseteq M\). The image \(e(T)\) is finite, and
  it \(R'\)-spans \(M'\): for \(m' \in M'\) write \(m' = e(m)\) with
  \(m = \sum_i r_i t_i\) (\(t_i \in T\)); then
  \(m' = e(\sum_i r_i t_i) = \sum_i \sigma(r_i)\, e(t_i)\) by semilinearity, an
  \(R'\)-combination of \(e(T)\). Hence the top submodule of \(M'\) is finitely
  generated, i.e. \(M'\) is a finite \(R'\)-module.
\end{proof}

\begin{lemma}\leanok
[G1 base case — finite affine sections from a finite generating family]
  \label{lem:gf_qcoh_finite_sections_of_genSections}
  \lean{AlgebraicGeometry.gf_qcoh_finite_sections_of_genSections}
  \uses{lem:gf_qcoh_sections_free_epi, lem:qcoh_affine_isIso_fromTildeΓ,
        lem:module_finite_of_ringEquiv_semilinear}
  Let \(X\) be a scheme, \(\F : X.\mathrm{Modules}\) a quasi-coherent sheaf of modules, and
  \(D \subseteq X\) an affine open. Suppose the restriction
  \(\F|_D = (\mathrm{pullback}\,D.\iota)^{*}\F\) carries a finite generating family
  \(\sigma\), with index set \(I = \sigma.I\). Then the section module
  \(\Gamma(\F, D)\) is a finite module over \(\Gamma(X, D)\). This is the affine base
  case of \(G1\): once the generators are in hand on an affine open, finiteness of
  sections is purely an affine (Spec-side) statement. (Only quasi-coherence of \(\F\) is
  needed — the finite-generation datum is carried by \(\sigma\) directly, so no separate
  finite-type hypothesis on \(\F\) is required.)
\end{lemma}

\begin{proof}
  \uses{lem:gf_qcoh_sections_free_epi, lem:qcoh_affine_isIso_fromTildeΓ,
        lem:module_finite_of_ringEquiv_semilinear, lem:qcoh_section_localization_basicOpen}
        \leanok
  The content is the transport of the finite generating datum
  from the category of sheaves of modules over the affine open \(D\) to the affine
  scheme \(\Spec \Gamma(X, D)\), where the affine base case applies. Write
  \(R = \Gamma(X, D)\) and use the canonical identification
  \(D \cong \Spec R\) (\(\mathrm{isoSpec}\)). The argument has three steps.

  \emph{(a) Quasi-coherence on \(\Spec R\).} Transport \(\F|_D\) along
  \(\mathrm{isoSpec}^{-1}\) to a sheaf of modules \(\F'\) on \(\Spec R\). Finite type of
  \(\F\) is inherited by the restriction, and quasi-coherence of \(\F'\) is the assertion
  that its counit \(\F'.\mathrm{fromTildeΓ}\) is an isomorphism; this is the affine
  quasi-coherence criterion \cref{lem:qcoh_affine_isIso_fromTildeΓ} fed by the basic-open
  localization datum \cref{lem:qcoh_section_localization_basicOpen}.

  \emph{(b) A finite free epimorphism over \(\Spec R\).} The generating family \(\sigma\)
  furnishes an epimorphism \(\OO_D^{\oplus I} \twoheadrightarrow \F|_D\). Transporting along
  \(\mathrm{isoSpec}^{-1}\) and using that the pullback of the free module along an
  isomorphism is again free (\((\mathrm{pullback}\,\mathrm{isoSpec}^{-1})^{*}\OO_D^{\oplus I}
  \cong \OO_{\Spec R}^{\oplus I}\)), together with the identification of the free sheaf with
  the associated sheaf \(\widetilde{N}\) of \(N = R^{\oplus I}\) (the associated-sheaf
  functor preserves coproducts and \(\widetilde{R} \cong \OO_{\Spec R}\)), gives an
  epimorphism \(\widetilde{N} \twoheadrightarrow \F'\) with \(N = R^{\oplus I}\) a finite
  \(R\)-module (\(I\) finite).

  \emph{(c) Identification of section modules.} The affine global-sections functor sends
  \(\F'\) to \(\Gamma(\F', \top)\), which the transport identifies with \(\Gamma(\F, D)\)
  as \(R = \Gamma(X, D)\)-modules.

  Now the affine base case \cref{lem:gf_qcoh_sections_free_epi}, applied to the finite-free
  quasi-coherent epimorphism \(\widetilde{N} \twoheadrightarrow \F'\) over \(\Spec R\),
  shows \(\Gamma(\F', \top)\) is a finite \(R\)-module; by (c) this is \(\Gamma(\F, D)\)
  finite over \(\Gamma(X, D)\).
\end{proof}

\begin{lemma}\leanok
[G1 — finite-type quasi-coherent sheaf has finite affine sections]
  \label{lem:gf_qcoh_fintype_finite_sections}
  \lean{AlgebraicGeometry.gf_qcoh_fintype_finite_sections}
  \uses{lem:gf_finite_sections_of_basicOpen_finite_cover,
        lem:gf_finiteType_affine_finite_cover_generated,
        lem:gf_qcoh_sections_free_epi,
        lem:gf_qcoh_finite_sections_of_genSections}
  % SOURCE: [Stacks Project], "Properties of Schemes", Tag 01PB
  % (Lemma "Characterizing modules of finite type", \S "Characterizing modules of
  % finite type and finite presentation")
  % (read from references/stacks-properties.tex, L2092-L2098).
  % SOURCE QUOTE:
  % "Let $X = \Spec(R)$ be an affine scheme.
  %  The quasi-coherent sheaf of $\mathcal{O}_X$-modules
  %  $\widetilde M$ is a finite type $\mathcal{O}_X$-module
  %  if and only if $M$ is a finite $R$-module."
  \textit{Source: [Stacks Project], "Properties of Schemes", Tag 01PB.}
  Let \(X\) be a scheme and \(\F : X.\mathrm{Modules}\) a quasi-coherent sheaf of
  modules with \(\F\) of finite type. Then for every affine open \(W \subseteq X\), the
  section module \(\Gamma(\F, W)\) is a finite module over \(\Gamma(X, W)\).
\end{lemma}

\begin{proof}
  \uses{lem:gf_finite_sections_of_basicOpen_finite_cover,
        lem:gf_finiteType_affine_finite_cover_generated,
        lem:gf_qcoh_sections_free_epi,
        lem:gf_qcoh_finite_sections_of_genSections}
        \leanok
  % SOURCE QUOTE PROOF: see the comment above the statement (Tag 01PB, L2092-L2098;
  % the Stacks proof reduces finite-type on a standard open cover to the algebra
  % gluing lemma for finite generation of $M$ from finite generation of each $M_{f_i}$).
  \textit{Source: [Stacks Project], "Properties of Schemes", Tag 01PB.}
  Fix an affine open \(W \subseteq X\). Cover extraction
  \cref{lem:gf_finiteType_affine_finite_cover_generated} gives a finite
  \(t \subseteq \Gamma(X, W)\) spanning the unit ideal, with each \(D(g)\) (\(g \in t\))
  affine and \(\F|_{D(g)}\) globally generated by finitely many sections, equivalently a
  finite-free epimorphism \(\OO_{D(g)}^{\oplus I_g} \twoheadrightarrow \F|_{D(g)}\). The
  free-epi base case \cref{lem:gf_qcoh_sections_free_epi} applied on each affine \(D(g)\)
  makes \(\Gamma(\F, D(g))\) finite over \(\Gamma(X, D(g))\). The locality reduction
  \cref{lem:gf_finite_sections_of_basicOpen_finite_cover} then glues these per-member
  finiteness statements to give \(\Gamma(\F, W)\) finite over \(\Gamma(X, W)\).
\end{proof}

\begin{lemma}[Mathlib anchor — free \(\Rightarrow\) flat]
  \label{lem:mathlib_flat_of_free}
  \lean{Module.Flat.of_free}
  \mathlibok
  \textit{Provided by Mathlib (\texttt{Mathlib.RingTheory.Flat.Basic}).}
  Let \(R\) be a commutative ring and \(M\) a free \(R\)-module. Then \(M\) is flat
  over \(R\).
\end{lemma}

\begin{lemma}[Mathlib anchor — localization preserves flatness]
  \label{lem:mathlib_flat_localization_preserves}
  \lean{Module.Flat.of_isLocalizedModule}
  \mathlibok
  \textit{Provided by Mathlib (\texttt{Mathlib.RingTheory.Flat.Stability}).}
  Let \(M\) be a flat \(R\)-module, let \(S \subseteq R\) be a multiplicative subset,
  let \(R_S\) be the localization of \(R\) at \(S\), and let \(M_S\) be a localization
  of \(M\) at \(S\) (an \(R_S\)-module via the scalar tower \(R \to R_S\)). Then
  \(M_S\) is flat over \(R_S\).
\end{lemma}

\begin{lemma}[Mathlib anchor — a localization is a flat module]
  \label{lem:mathlib_localization_flat}
  \lean{IsLocalization.flat}
  \mathlibok
  \textit{Provided by Mathlib (\texttt{Mathlib.RingTheory.Flat.Localization}).}
  Let \(R\) be a commutative ring, \(S \subseteq R\) a multiplicative subset, and
  \(R_S\) the localization of \(R\) at \(S\). Then \(R_S\) is flat as an
  \(R\)-module.
\end{lemma}

\begin{lemma}[Mathlib anchor — flatness is local at maximal ideals of the base]
  \label{lem:mathlib_flat_of_localized_maximal}
  \lean{Module.flat_of_localized_maximal}
  \mathlibok
  \textit{Provided by Mathlib (\texttt{Mathlib.RingTheory.Flat.Localization}).}
  Let \(R\) be a commutative ring and \(M\) an \(R\)-module. If for every maximal
  ideal \(\mathfrak p \subset R\) the localization
  \(M_{\mathfrak p} = \mathrm{LocalizedModule}\,\mathfrak p^c\, M\) is flat over \(R\),
  then \(M\) is flat over \(R\).
\end{lemma}

\begin{lemma}[Mathlib anchor — flatness over a base is local on the source]
  \label{lem:mathlib_flat_of_isLocalized_maximal}
  \lean{Module.flat_of_isLocalized_maximal}
  \mathlibok
  \textit{Provided by Mathlib (\texttt{Mathlib.RingTheory.Flat.Localization}).}
  Let \(S\) be an \(R\)-algebra and \(M\) an \(S\)-module, regarded as an
  \(R\)-module through the scalar tower \(R \to S \to M\). Suppose that for every
  maximal ideal \(\mathfrak q \subset S\) the localized module \(M_{\mathfrak q}\)
  (the localization of \(M\) at the prime \(\mathfrak q\) of \(S\), an
  \(S_{\mathfrak q}\)-module and hence an \(R\)-module) is flat over \(R\). Then
  \(M\) is flat over \(R\). \emph{(Flatness over the base ring \(R\) is detected at
  the maximal ideals of the source ring \(S\).)}
\end{lemma}

\begin{lemma}[Mathlib anchor — transitivity of flatness]
  \label{lem:mathlib_flat_trans}
  \lean{Module.Flat.trans}
  \mathlibok
  \textit{Provided by Mathlib (\texttt{Mathlib.RingTheory.Flat.Stability}).}
  Let \(R \to S\) be a ring homomorphism and \(M\) an \(S\)-module, viewed as an
  \(R\)-module through the scalar tower. If \(S\) is flat over \(R\) and \(M\) is
  flat over \(S\), then \(M\) is flat over \(R\).
\end{lemma}

\begin{lemma}[Mathlib anchor — flatness over a base is local on a spanning family of the source]
  \label{lem:mathlib_flat_of_isLocalized_span}
  \lean{Module.flat_of_isLocalized_span}
  \mathlibok
  \textit{Provided by Mathlib (\texttt{Mathlib.RingTheory.Flat.Localization}).}
  Let \(R\) be a commutative ring, \(S\) an \(R\)-algebra, and \(M\) an \(S\)-module,
  regarded as an \(R\)-module through the scalar tower \(R \to S \to M\). Let
  \(s \subseteq S\) be a subset whose ideal span is all of \(S\)
  (\(\mathrm{Ideal.span}\,s = \top\)), and for each \(r \in s\) let \(M_r\) be an
  away-localization of \(M\) at \(r\) (an \(S_r\)-module, hence an \(R\)-module via the
  tower). If \(M_r\) is flat over \(R\) for every \(r \in s\), then \(M\) is flat over
  \(R\). \emph{(Flatness over the base ring \(R\) descends from a spanning family of
  principal localizations of the \emph{source} ring \(S\).)}
\end{lemma}

\begin{lemma}

\leanok
[Flatness transports up a localization tower of the base]
  \label{lem:gf_stalk_flat_localBase}
  \lean{AlgebraicGeometry.gf_stalk_flat_localBase}
  \uses{lem:mathlib_localization_flat, lem:mathlib_flat_trans}
  Let \(R \to R'\) exhibit \(R'\) as the localization of \(R\) at a submonoid
  \(S \subseteq R\), and let \(N\) be an \(R'\)-module, regarded as an \(R\)-module
  through the scalar tower \(R \to R' \to N\). If \(N\) is flat over \(R'\), then \(N\)
  is flat over \(R\). \emph{(Geometric use: with \(R = \OO_{S,x}\),
  \(R' = \OO_{S,p(y)}\) a localization of \(R\) at a generization, and \(N = \F_y\), this
  upgrades flatness of the stalk \(\F_y\) over the local base \(\OO_{S,p(y)}\) to flatness
  over \(\OO_{S,x}\).)}
\end{lemma}

\begin{proof}
  \uses{lem:mathlib_localization_flat, lem:mathlib_flat_trans}
  \leanok
  The localization \(R'\) is flat over \(R\) (\cref{lem:mathlib_localization_flat}). Since
  \(N\) is flat over \(R'\) by hypothesis, transitivity of flatness along the scalar tower
  \(R \to R' \to N\) (\cref{lem:mathlib_flat_trans}) yields that \(N\) is flat over \(R\).
\end{proof}

\begin{lemma}[Mathlib anchor — a localized module of a free module is free]
  \label{lem:mathlib_free_of_isLocalizedModule}
  \lean{Module.free_of_isLocalizedModule}
  \mathlibok
  \textit{Provided by Mathlib (\texttt{Mathlib.RingTheory.LocalProperties.Projective}).}
  Let \(R\) be a commutative ring, \(S \subseteq R\) a submonoid, \(R_S\) the localization
  of \(R\) at \(S\), and \(M\) a free \(R\)-module. Let \(f : M \to M_S\) be an \(R\)-linear
  map exhibiting \(M_S\) as the localization of \(M\) at \(S\) (an \(R_S\)-module through the
  scalar tower \(R \to R_S\)). Then \(M_S\) is free over \(R_S\).
\end{lemma}

\begin{lemma}\leanok
[G3.1 — free patch implies flat patch]
  \label{lem:gf_patch_free_imp_flat}
  \lean{AlgebraicGeometry.gf_patch_free_imp_flat}
  \uses{lem:mathlib_flat_of_free}
  In the situation of the generic-flatness assembly \cref{thm:generic_flatness}, for every
  \(j\) the localized section module \((M_j)_f\) is flat over \(A_f\).
\end{lemma}

\begin{proof}
  \uses{lem:mathlib_flat_of_free}
  \leanok
  By hypothesis \((M_j)_f\) is free over \(A_f\), and a free module over a
  commutative ring is flat (\cref{lem:mathlib_flat_of_free}). Hence \((M_j)_f\) is
  flat over \(A_f\).
\end{proof}

\begin{lemma}\leanok
[G3.3 — flatness over a local base is local on the source affine]
  \label{lem:gf_flat_base_local_on_source}
  \lean{AlgebraicGeometry.gf_flat_base_local_on_source}
  \uses{lem:mathlib_flat_of_isLocalized_maximal}
  Let \(R\) be a commutative ring, \(B\) an \(R\)-algebra, and \(N\) a \(B\)-module,
  regarded as an \(R\)-module through the scalar tower \(R \to B \to N\). Suppose
  that for every maximal ideal \(\mathfrak q \subset B\) the localization
  \(N_{\mathfrak q}\) is flat over \(R\). Then \(N\) is flat over \(R\).
\end{lemma}

\begin{proof}
  \uses{lem:mathlib_flat_of_isLocalized_maximal}
  \leanok
  This is the source-locality criterion for flatness over a fixed base ring,
  instantiated with the canonical localizations \(N_{\mathfrak q}\) of \(N\) at the
  maximal ideals \(\mathfrak q\) of \(B\): each \(N_{\mathfrak q}\) is a localized
  module of \(N\) at the prime \(\mathfrak q\) of the source ring \(B\), and the
  hypothesis supplies the required flatness of every \(N_{\mathfrak q}\) over \(R\).
  \Cref{lem:mathlib_flat_of_isLocalized_maximal} then yields that \(N\) is flat over
  \(R\). \emph{(Note: unlike \cref{lem:mathlib_flat_of_localized_maximal}, which
  localizes over the same ring it tests flatness over, here \(N\) is localized at the
  primes of the source ring \(B\) while flatness is tested over the distinct base
  ring \(R\).)}
\end{proof}

\begin{lemma}\leanok
[B1.0 — localization commutes with base tensor]
  \label{lem:gf_localizedModule_baseChange_tensor_comm}
  \lean{AlgebraicGeometry.gf_localizedModule_baseChange_tensor_comm}
  Let \(R \to B\) be a ring map, \(T \subseteq B\) a submonoid, \(N\) a \(B\)-module,
  and \(K\) an \(R\)-module. There is an isomorphism, natural in \(K\),
  \[
    \mathrm{LocalizedModule}\,T\,(N \otimes_R K)
      \;\cong\; (\mathrm{LocalizedModule}\,T\,N) \otimes_R K .
  \]
\end{lemma}
\begin{proof}

\leanok
  Localization is base change: \(\mathrm{LocalizedModule}\,T\,M \cong T^{-1}B \otimes_B M\)
  naturally in the \(B\)-module \(M\). Applying this to \(M = N \otimes_R K\) and using
  associativity/commutativity of the tensor product over the tower
  \(R \to B \to T^{-1}B\) gives
  \(T^{-1}B \otimes_B (N \otimes_R K) \cong (T^{-1}B \otimes_B N) \otimes_R K
  \cong (\mathrm{LocalizedModule}\,T\,N) \otimes_R K\); naturality in \(K\) is inherited
  from the tensor functor.
\end{proof}

\begin{lemma}\leanok
[B1 — localizing the source ring preserves flatness over the base]
  \label{lem:gf_flat_localizedModule_sameBase}
  \lean{AlgebraicGeometry.gf_flat_localizedModule_sameBase}
  \uses{lem:gf_localizedModule_baseChange_tensor_comm}
  Let \(R \to B \to N\) be a scalar tower in which \(N\) is a \(B\)-module, and let
  \(T \subseteq B\) be a submonoid of the \emph{source} ring \(B\). View \(N\) as an
  \(R\)-module through the tower. If \(N\) is flat over \(R\), then the localized module
  \(\mathrm{LocalizedModule}\,T\,N\) is again flat over \(R\). \emph{(Note: the existing
  Mathlib results \cref{lem:mathlib_flat_localization_preserves} localize over a
  submonoid of the \emph{base} \(R\); here the submonoid lives in the source \(B\) while
  the base \(R\) is untouched, so a separate argument is required.)}
\end{lemma}

\begin{proof}
  \uses{lem:gf_localizedModule_baseChange_tensor_comm, lem:mathlib_flat_localization_preserves}
  \leanok
  Localization at \(T\) commutes with tensoring against an \(R\)-module: by
  \cref{lem:gf_localizedModule_baseChange_tensor_comm}, for every \(R\)-module \(K\)
  there is a natural isomorphism
  \[
    \mathrm{LocalizedModule}\,T\,(N \otimes_R K)
      \;\cong\; (\mathrm{LocalizedModule}\,T\,N) \otimes_R K,
  \]
  natural in \(K\) (compatible with the maps induced by any \(R\)-linear \(K \to K'\)).
  Flatness of \(N\) over \(R\) means that for every injection \(K \hookrightarrow K'\) of
  \(R\)-modules the induced map \(N \otimes_R K \to N \otimes_R K'\) is injective.
  Localization at the submonoid \(T\) is an exact functor, so it preserves this
  injectivity; transporting along the natural isomorphism turns the result into
  injectivity of
  \((\mathrm{LocalizedModule}\,T\,N) \otimes_R K \to (\mathrm{LocalizedModule}\,T\,N) \otimes_R K'\).
  As this holds for every injection of \(R\)-modules, \(\mathrm{LocalizedModule}\,T\,N\)
  is flat over \(R\).
\end{proof}

\begin{lemma}

\leanok
[B1\(^\prime\) — model-free form of B1 for an arbitrary localization model]
  \label{lem:gf_flat_isLocalizedModule_sameBase}
  \lean{AlgebraicGeometry.gf_flat_isLocalizedModule_sameBase}
  \uses{lem:gf_flat_localizedModule_sameBase}
  Let \(R \to B \to N\) be a scalar tower with \(N\) a \(B\)-module, let \(T \subseteq B\) be a
  submonoid of the \emph{source} ring \(B\), and let \(\varphi : N \to N'\) be a \(B\)-linear map
  exhibiting \(N'\) as a localization of \(N\) at \(T\) — that is, \(N'\) is \emph{any}
  localized-module model of \(N\) at \(T\), rather than the canonical
  \(\mathrm{LocalizedModule}\,T\,N\). Give \(N'\) the \(R\)-module structure induced through the
  tower. If \(N\) is flat over \(R\), then \(N'\) is flat over \(R\). \emph{(This is the
  consumer-facing, model-free variant of \cref{lem:gf_flat_localizedModule_sameBase}, applied in
  \cref{thm:generic_flatness} to a geometric section module \(\Gamma(\F, D)\) presented as a
  localization.)}
\end{lemma}

\begin{proof}
  \uses{lem:gf_flat_localizedModule_sameBase}
  \leanok
  By \cref{lem:gf_flat_localizedModule_sameBase} the canonical localized module
  \(\mathrm{LocalizedModule}\,T\,N\) is flat over \(R\). Uniqueness of localizations supplies the
  canonical comparison isomorphism \(N' \cong \mathrm{LocalizedModule}\,T\,N\) between any two
  localizations of \(N\) at \(T\); restricting its scalars to \(R\) gives an \(R\)-linear
  isomorphism. Flatness over \(R\) transports across this isomorphism, so \(N'\) is flat over
  \(R\).
\end{proof}

\begin{lemma}

\leanok
[Flatness transports across a semilinear ring isomorphism]
  \label{lem:flat_of_ringEquiv_semilinear}
  \lean{AlgebraicGeometry.flat_of_ringEquiv_semilinear}
  \uses{lem:module_free_of_ringEquiv}
  Let \(R, R'\) be commutative rings, \(M\) an \(R\)-module and \(M'\) an \(R'\)-module. Let
  \(e : R \xrightarrow{\sim} R'\) be a ring isomorphism and \(l : M \xrightarrow{\sim} M'\) an
  additive isomorphism that is \(e\)-semilinear, i.e.\ \(l(r \cdot x) = e(r) \cdot l(x)\) for all
  \(r \in R\) and \(x \in M\). If \(M\) is flat over \(R\), then \(M'\) is flat over \(R'\). \emph{(This
  is the flatness analogue of the freeness-transport \cref{lem:module_free_of_ringEquiv}.)}
\end{lemma}

\begin{proof}
  \uses{lem:module_free_of_ringEquiv}
  \leanok
  Regard \(R'\) as an \(R\)-algebra through \(e\) and \(M'\) as an \(R\)-module by restriction of
  scalars along \(e\); then \(l\) is \(R\)-linear, since \(l(r \cdot x) = e(r) \cdot l(x)\) is exactly
  \(R\)-linearity for the restricted action. This presents \(M'\) as the base change
  \(R' \otimes_R M\): the \(R'\)-linear lift of \(l\) (sending \(r' \otimes m \mapsto r' \cdot l(m)\))
  has two-sided inverse \(y \mapsto 1 \otimes l^{-1}(y)\), using the \(e^{-1}\)-semilinearity of
  \(l^{-1}\), so \(l\) exhibits \(M'\) as \(R' \otimes_R M\). The base change of a flat module along a
  ring map is flat over the new base ring; hence \(M'\) is flat over \(R'\).
\end{proof}

\begin{lemma}

\leanok
[Flatness of a localization is independent of the chosen models]
  \label{lem:flat_localization_models}
  \lean{AlgebraicGeometry.flat_localization_models}
  \uses{lem:flat_of_ringEquiv_semilinear}
  Let \(M\) be a module over a commutative ring \(A\) and \(S \subseteq A\) a submonoid. Let
  \(M_S, M_S'\) be two localizations of \(M\) at \(S\), taken over the respective localizations
  \(R_S, R_S'\) of \(A\) at \(S\). If \(M_S\) is flat over \(R_S\), then \(M_S'\) is flat over
  \(R_S'\).
\end{lemma}

\begin{proof}
  \uses{lem:flat_of_ringEquiv_semilinear}
  \leanok
  By uniqueness of localizations, the rings \(R_S\) and \(R_S'\) are canonically isomorphic as
  \(A\)-algebras, giving a ring isomorphism \(e : R_S \xrightarrow{\sim} R_S'\); likewise the
  modules \(M_S\) and \(M_S'\) are canonically isomorphic as \(A\)-modules, giving an additive
  isomorphism \(l : M_S \xrightarrow{\sim} M_S'\). These comparison maps are \(e\)-semilinear: for a
  scalar \(c \in R_S\), write \(c\) as a fraction with denominator \(s \in S\); since \(s\) acts
  invertibly on both \(M_S\) and \(M_S'\), the identity \(l(c \cdot x) = e(c) \cdot l(x)\) follows
  by clearing \(s\) and using that \(e\) and \(l\) are \(A\)-linear on the resulting \(A\)-scalars.
  The semilinear transport \cref{lem:flat_of_ringEquiv_semilinear} then carries flatness of \(M_S\)
  over \(R_S\) to flatness of \(M_S'\) over \(R_S'\).
\end{proof}

\begin{lemma}

\leanok
[Scalar-tower descent of a powers localization]
  \label{lem:isLocalizedModule_powers_restrictScalars}
  \lean{AlgebraicGeometry.isLocalizedModule_powers_restrictScalars}
  Let \(A \to B\) be a ring map, let \(M, N\) be \(B\)-modules with the \(A\)-module structures
  induced through the tower, and let \(f \in A\). If a \(B\)-linear map \(\varphi : M \to N\)
  exhibits \(N\) as the localization of \(M\) at the powers of the image
  \(\mathrm{algebraMap}_{A \to B}(f) \in B\), then the scalar restriction
  \(\varphi|_A : M \to N\) exhibits \(N\) as the localization of \(M\), viewed as an \(A\)-module,
  at the powers of \(f\).
\end{lemma}

\begin{proof}

\leanok
  The \(A\)-module structures on \(M\) and \(N\) factor through the \(B\)-action via the scalar
  tower, so for each \(k \in \mathbb{N}\) the action of \(f^k \in A\) coincides with the action of
  \(\mathrm{algebraMap}_{A \to B}(f)^k = \mathrm{algebraMap}_{A \to B}(f^k) \in B\). The three
  defining conditions of a localized module — invertibility of the action of each power of \(f\),
  surjectivity onto \(N\) by fractions with such denominators, and the equality criterion — for the
  submonoid \(\mathrm{powers}\,f \subseteq A\) therefore agree element-by-element with the
  corresponding conditions for \(\mathrm{powers}\,(\mathrm{algebraMap}_{A \to B}(f)) \subseteq B\),
  which hold by hypothesis. Hence \(\varphi|_A\) exhibits \(N\) as the localization of \(M\) at
  \(\mathrm{powers}\,f\).
\end{proof}

\begin{lemma}[Mathlib anchor — basic open under restriction]
  \label{lem:mathlib_scheme_basicOpen_res}
  \lean{AlgebraicGeometry.Scheme.basicOpen_res}
  \mathlibok
  \textit{Provided by Mathlib (\texttt{Mathlib.AlgebraicGeometry.Scheme}).}
  Let \(X\) be a scheme, \(V \le U\) opens of \(X\), and \(s \in \Gamma(X, U)\). Restricting
  \(s\) to \(V\) meets its basic open with \(V\):
  \[
    X.\mathrm{basicOpen}\,(s|_V) \;=\; V \sqcap X.\mathrm{basicOpen}\,s .
  \]
\end{lemma}

\begin{lemma}\leanok
[B2.1 — cross-chart basic-open identity]
  \label{lem:gf_crossChart_basicOpen_eq}
  \lean{AlgebraicGeometry.gf_crossChart_basicOpen_eq}
  \uses{lem:mathlib_scheme_basicOpen_res}
  Let \(X\) be a scheme, \(W, W_j \subseteq X\) two opens with overlap
  \(O = W \sqcap W_j\), and let \(g \in \Gamma(X, W)\), \(\bar g \in \Gamma(X, W_j)\) be sections
  that restrict to the \emph{same} section over \(O\), i.e.\ \(g|_O = \bar g|_O\) in
  \(\Gamma(X, O)\). Suppose both basic opens lie in the overlap,
  \(X.\mathrm{basicOpen}\,g \le O\) and \(X.\mathrm{basicOpen}\,\bar g \le O\). Then they agree as
  opens of \(X\):
  \[
    X.\mathrm{basicOpen}\,g \;=\; X.\mathrm{basicOpen}\,\bar g .
  \]
  In particular the open \(D(g)\) is computed identically whether read in the chart \(W\) or in
  the chart \(W_j\), and \(D(g) = D(\bar g) \le W_j\) (no module content is involved).
\end{lemma}

\begin{proof}
  \uses{lem:mathlib_scheme_basicOpen_res}
  \leanok
  Basic opens are compatible with restriction (\cref{lem:mathlib_scheme_basicOpen_res}). For the
  inclusion \(O \le W\), \(X.\mathrm{basicOpen}\,(g|_O) = O \sqcap X.\mathrm{basicOpen}\,g\), and
  since \(X.\mathrm{basicOpen}\,g \le O\) the meet collapses to
  \(X.\mathrm{basicOpen}\,(g|_O) = X.\mathrm{basicOpen}\,g\). Symmetrically, for \(O \le W_j\),
  \(X.\mathrm{basicOpen}\,(\bar g|_O) = O \sqcap X.\mathrm{basicOpen}\,\bar g
  = X.\mathrm{basicOpen}\,\bar g\). The hypothesis \(g|_O = \bar g|_O\) makes the two left-hand
  sides equal, so \(X.\mathrm{basicOpen}\,g = X.\mathrm{basicOpen}\,\bar g\).
\end{proof}

\begin{lemma}\leanok
[B2.2 — two-leg section-localization transport]
  \label{lem:gf_section_localization_twoleg}
  \lean{AlgebraicGeometry.gf_section_localization_twoleg}
  \uses{lem:gf_crossChart_basicOpen_eq, lem:qcoh_section_localization_basicOpen,
        lem:gf_qcoh_fintype_finite_sections}
  In the situation of the generic-flatness assembly \cref{thm:generic_flatness}, fix an affine open
  \(W \le p^{-1}(U)\), a patch \(W_j = \Spec B_j\) with \(M_j = \Gamma(\F, W_j)\), and a matched
  pair \(g \in \Gamma(X, W)\), \(\bar g \in B_j = \Gamma(X, W_j)\) cutting out the same basic open
  \(D := X.\mathrm{basicOpen}\,g = X.\mathrm{basicOpen}\,\bar g \le W \sqcap W_j\) (as produced by
  \cref{lem:gf_crossChart_basicOpen_eq}). Then the single group \(\Gamma(\F, D)\) carries two
  localized-module structures:
  \begin{enumerate}
    \item the section restriction \(\Gamma(\F, W) \to \Gamma(\F, D)\) exhibits \(\Gamma(\F, D)\) as
      \((\mathrm{powers}\,g)^{-1}\,\Gamma(\F, W)\); and
    \item the section restriction \(M_j \to \Gamma(\F, D)\) exhibits \(\Gamma(\F, D)\) as
      \((\mathrm{powers}\,\bar g)^{-1}\,M_j = (M_j)_{\bar g}\).
  \end{enumerate}
  Consequently \((\mathrm{powers}\,g)^{-1}\,\Gamma(\F, W) \cong (M_j)_{\bar g}\) canonically, by
  uniqueness of localized modules.
\end{lemma}

\begin{proof}
  \uses{lem:gf_crossChart_basicOpen_eq, lem:qcoh_section_localization_basicOpen,
        lem:gf_qcoh_fintype_finite_sections}
        \leanok
  By \cref{lem:gf_crossChart_basicOpen_eq} the two recipes for the open agree,
  \(D = X.\mathrm{basicOpen}\,g = X.\mathrm{basicOpen}\,\bar g\), so \(\Gamma(\F, D)\) is one and the
  same group whether \(D\) is read in chart \(W\) or in chart \(W_j\). Applying the quasi-coherent
  section-localization datum \cref{lem:qcoh_section_localization_basicOpen} with ambient affine \(W\)
  and element \(g\), the restriction \(\Gamma(\F, W) \to \Gamma(\F, D)\) is the localization at
  \(\mathrm{powers}\,g\), giving (1). Applying the same datum with ambient affine \(W_j\) and element
  \(\bar g\) (legitimate since \(D = X.\mathrm{basicOpen}\,\bar g \le W_j\)), the restriction
  \(M_j \to \Gamma(\F, D)\) is the localization at \(\mathrm{powers}\,\bar g\), giving (2). Both maps
  localize their respective sources onto the same target \(\Gamma(\F, D)\); the universal property of
  localized modules (\(\mathrm{IsLocalizedModule.iso}\) uniqueness) supplies the canonical isomorphism
  \((\mathrm{powers}\,g)^{-1}\,\Gamma(\F, W) \cong (M_j)_{\bar g}\). Finiteness of \(M_j\) over \(B_j\)
  is recorded by \cref{lem:gf_qcoh_fintype_finite_sections}.
\end{proof}

\begin{lemma}\leanok
[B2.3 — base comparison: \(\Gamma(S,U)\) localizes \(A_f\)]
  \label{lem:gf_base_localization_comparison}
  \lean{AlgebraicGeometry.gf_base_localization_comparison}
  \uses{lem:isLocalization_basicOpen_mathlib}
  % NOTE (iter-054): the Lean proves the strictly WEAKER form `Module.Flat Γ(S,V) Γ(S,U)`
  % (flatness of the restriction algebra), NOT `IsLocalization` as the prose states. Flatness
  % is all the assembly consumes (localizations are flat ⟹ this suffices), and it is reached
  % directly via `Flat.flat_appLE (𝟙 S)`. Planner: downgrade the prose claim to flatness, or
  % strengthen the Lean to `IsLocalization`, to remove the name-vs-type mismatch.
  In the situation of the generic-flatness assembly \cref{thm:generic_flatness}, let
  \(U \le V = D(f) \subseteq U_0 = \Spec A\)
  be an affine open of \(S\). Then \(A_f = \Gamma(S, D(f))\) is the away-localization of \(A\) at
  \(f\), and the restriction \(A_f \to R := \Gamma(S, U)\) realises \(R\) as a localization of \(A_f\)
  (both are rings of fractions of the domain \(A\) inside its fraction field, \(R\) inverting at least
  the denominators already inverted in \(A_f\)).
\end{lemma}

\begin{proof}
  \uses{lem:isLocalization_basicOpen_mathlib}
  \leanok
  Since \(D(f)\) is a basic open of the affine \(U_0 = \Spec A\), the restriction
  \(A = \Gamma(S, U_0) \to \Gamma(S, D(f))\) exhibits \(A_f = \Gamma(S, D(f))\) as the
  away-localization of \(A\) at \(f\) (\cref{lem:isLocalization_basicOpen_mathlib}). The affine open
  \(U \le D(f) = \Spec A_f\) of the integral scheme \(S\) has section ring
  \(R = \Gamma(S, U) \subseteq \mathrm{Frac}(A)\) containing \(A_f\); the comparison
  \(A_f \to R\) is the restriction along \(U \le D(f)\) and exhibits \(R\) as a localization of
  \(A_f\) inside the fraction field of the domain \(A\).
\end{proof}

\subsection{Base descent along the open immersion \(U \hookrightarrow V\)}
\label{sec:gf_base_descent_isEpi}

The generic-flatness construction yields, on each patch, flatness of the section module
over the base ring \(\Gamma(S, V) = A_f\) of the witness open \(V = D(f)\). For an
\emph{arbitrary} affine open \(U \le V\), the base ring \(R = \Gamma(S, U)\) need not be a
localization of \(A_f\), so the localization-descent of
\cref{lem:gf_base_localization_comparison} does not transport flatness directly. The correct
mechanism is that the inclusion \(U \hookrightarrow V\) of affine opens is an \emph{open
immersion}, hence a monomorphism of schemes, so the restriction ring map
\(\Gamma(S, V) \to \Gamma(S, U)\) is a ring \emph{epimorphism}. For a ring epimorphism
\(A \to R\) the canonical base-change map \(R \otimes_A M \to M\) is an isomorphism, and
flatness descends along it. This replaces the earlier \texttt{IsBaseChange} hypothesis (once
flagged as absent from Mathlib) by the verified Mathlib epimorphism API.

\begin{definition}[Mathlib anchor — ring-epimorphism typeclass]
  \label{lem:mathlib_algebra_isEpi}
  \lean{Algebra.IsEpi}
  \mathlibok
  \textit{Provided by Mathlib (\texttt{Mathlib.Algebra.Algebra.Epi}).}
  For a commutative ring \(R\) and an \(R\)-algebra \(A\), the predicate
  \(\mathrm{Algebra.IsEpi}\,R\,A\) asserts that the structure map \(R \to A\) is an
  epimorphism in the category of commutative rings (equivalently, the multiplication map
  \(A \otimes_R A \to A\) is an isomorphism).
\end{definition}

\begin{lemma}[Mathlib anchor — base change along a ring epimorphism is the identity]
  \label{lem:mathlib_tensorProduct_lid_epi}
  \lean{TensorProduct.lid'}
  \mathlibok
  \uses{lem:mathlib_algebra_isEpi}
  \textit{Provided by Mathlib (\texttt{Mathlib.Algebra.Algebra.Epi}).}
  Let \(R \to A\) be a ring epimorphism (\(\mathrm{Algebra.IsEpi}\,R\,A\)) and \(M\) an
  \(A\)-module regarded as an \(R\)-module through the scalar tower \(R \to A \to M\). Then
  the canonical map exhibits the base change \(A \otimes_R M\) as \(M\) itself, an
  isomorphism of \(A\)-modules
  \[
    \mathrm{TensorProduct.lid'} : A \otimes_R M \;\xrightarrow{\ \cong\ }\; M .
  \]
\end{lemma}

\begin{lemma}[Mathlib anchor — categorical epi of a ring map is a ring epimorphism]
  \label{lem:mathlib_commRingCat_epi_iff_epi}
  \lean{CommRingCat.epi_iff_epi}
  \mathlibok
  \uses{lem:mathlib_algebra_isEpi}
  \textit{Provided by Mathlib (\texttt{Mathlib.Algebra.Category.Ring.Epi}).}
  For commutative rings \(R, S\) with an algebra structure \(R \to S\), the ring map
  \(R \to S\) is an epimorphism in \(\mathrm{CommRingCat}\) if and only if
  \(\mathrm{Algebra.IsEpi}\,R\,S\) holds:
  \[
    \mathrm{Epi}\,(\mathrm{CommRingCat.ofHom}\,(\mathrm{algebraMap}\,R\,S))
      \;\Longleftrightarrow\; \mathrm{Algebra.IsEpi}\,R\,S .
  \]
\end{lemma}

\begin{lemma}[Mathlib anchor — \(\Spec\) is fully faithful]
  \label{lem:mathlib_spec_fullyFaithful}
  \lean{AlgebraicGeometry.Spec.fullyFaithful}
  \mathlibok
  \textit{Provided by Mathlib (\texttt{Mathlib.AlgebraicGeometry.Spec}).}
  The functor \(\Spec : \mathrm{CommRingCat}^{\mathrm{op}} \to \mathrm{Scheme}\) is fully
  faithful. In particular a monomorphism \(\Spec S \to \Spec R\) of affine schemes
  corresponds to an epimorphism \(R \to S\) of rings.
\end{lemma}

\begin{lemma}[Mathlib anchor — an open immersion is a monomorphism]
  \label{lem:mathlib_isOpenImmersion_mono}
  \lean{AlgebraicGeometry.IsOpenImmersion.mono}
  \mathlibok
  \textit{Provided by Mathlib (\texttt{Mathlib.AlgebraicGeometry.Morphisms.OpenImmersion}).}
  Every open immersion \(f : X \to Z\) of schemes is a monomorphism in the category of
  schemes.
\end{lemma}

\begin{lemma}[Mathlib anchor — an affine open is its own \(\Spec\)]
  \label{lem:mathlib_isAffineOpen_isoSpec}
  \lean{AlgebraicGeometry.IsAffineOpen.isoSpec}
  \mathlibok
  \textit{Provided by Mathlib (\texttt{Mathlib.AlgebraicGeometry.AffineScheme}).}
  For an affine open \(U\) of a scheme \(X\) there is a canonical isomorphism of schemes
  \(U \cong \Spec \Gamma(X, U)\), identifying the open subscheme \(U\) with the spectrum of
  its ring of sections.
\end{lemma}

\begin{lemma}[Mathlib anchor — base change of a flat module is flat]
  \label{lem:mathlib_baseChange_flat}
  \lean{Module.Flat.baseChange}
  \mathlibok
  \textit{Provided by Mathlib (\texttt{Mathlib.RingTheory.Flat.Basic}).}
  Let \(A \to R\) be a ring map and \(M\) a flat \(A\)-module. Then the base change
  \(R \otimes_A M\) is flat over \(R\).
\end{lemma}

\begin{lemma}[Mathlib anchor — flatness transports across a linear equivalence]
  \label{lem:mathlib_flat_of_linearEquiv}
  \lean{Module.Flat.of_linearEquiv}
  \mathlibok
  \textit{Provided by Mathlib (\texttt{Mathlib.RingTheory.Flat.Basic}).}
  Let \(R\) be a commutative ring and \(M \cong N\) an isomorphism of \(R\)-modules. If
  \(M\) is flat over \(R\), then \(N\) is flat over \(R\).
\end{lemma}

\begin{lemma}[Mathlib anchor — the target of a base-change identification is flat]
  \label{lem:mathlib_flat_isBaseChange}
  \lean{Module.Flat.isBaseChange}
  \mathlibok
  \textit{Provided by Mathlib (\texttt{Mathlib.RingTheory.Flat.Stability}).}
  Let \(A \to R\) be a ring map, \(M\) a flat \(A\)-module, and \(N\) an \(R\)-module with
  an \(A\)-linear map \(f : M \to N\) exhibiting \(N\) as the base change of \(M\) along
  \(A \to R\) (\(\mathrm{IsBaseChange}\,R\,f\)). Then \(N\) is flat over \(R\).
\end{lemma}

\begin{lemma}\leanok
[Restriction between affine opens is a ring epimorphism]
  \label{lem:gf_openImmersion_isEpi}
  \lean{AlgebraicGeometry.gf_isEpi_restrict_of_affine_le}
  \uses{lem:mathlib_spec_fullyFaithful, lem:mathlib_isOpenImmersion_mono,
        lem:mathlib_isAffineOpen_isoSpec, lem:mathlib_commRingCat_epi_iff_epi,
        lem:mathlib_algebra_isEpi}
  Let \(U \le V\) be affine opens of a scheme \(S\). Then the structure-sheaf restriction
  \(\Gamma(S, V) \to \Gamma(S, U)\) is a ring epimorphism, i.e.\
  \(\mathrm{Algebra.IsEpi}\,\Gamma(S, V)\,\Gamma(S, U)\).
\end{lemma}

\begin{proof}
  \uses{lem:mathlib_spec_fullyFaithful, lem:mathlib_isOpenImmersion_mono,
        lem:mathlib_isAffineOpen_isoSpec, lem:mathlib_commRingCat_epi_iff_epi}
        \leanok
  The inclusion \(\iota : U \hookrightarrow V\) of opens of \(S\) is an open immersion of
  schemes, hence a monomorphism (\cref{lem:mathlib_isOpenImmersion_mono}). Both \(U\) and
  \(V\) are affine, so under the canonical identifications \(U \cong \Spec \Gamma(S, U)\)
  and \(V \cong \Spec \Gamma(S, V)\) (\cref{lem:mathlib_isAffineOpen_isoSpec}) the
  monomorphism \(\iota\) is \(\Spec\) of the restriction ring map
  \(\Gamma(S, V) \to \Gamma(S, U)\). Since \(\Spec\) is fully faithful
  (\cref{lem:mathlib_spec_fullyFaithful}), a monomorphism between affine schemes
  corresponds to an \emph{epimorphism} of section rings; thus
  \(\Gamma(S, V) \to \Gamma(S, U)\) is an epimorphism in \(\mathrm{CommRingCat}\). Finally
  \cref{lem:mathlib_commRingCat_epi_iff_epi} repackages this categorical epimorphism as the
  ring-epimorphism property \(\mathrm{Algebra.IsEpi}\,\Gamma(S, V)\,\Gamma(S, U)\).
\end{proof}

\begin{lemma}\leanok
[Flat descent along a ring epimorphism]
  \label{lem:gf_flat_descend_isEpi}
  \lean{AlgebraicGeometry.gf_flat_of_isEpi}
  \uses{lem:mathlib_algebra_isEpi, lem:mathlib_tensorProduct_lid_epi,
        lem:mathlib_baseChange_flat, lem:mathlib_flat_of_linearEquiv}
  Let \(\Gamma(S, V) \to \Gamma(S, U) \to M\) be a scalar tower in which \(M\) is a
  \(\Gamma(S, U)\)-module, with \(\mathrm{Algebra.IsEpi}\,\Gamma(S, V)\,\Gamma(S, U)\).
  If \(M\) is flat over \(\Gamma(S, V)\), then \(M\) is flat over \(\Gamma(S, U)\).
\end{lemma}

\begin{proof}
  \uses{lem:mathlib_tensorProduct_lid_epi, lem:mathlib_baseChange_flat,
        lem:mathlib_flat_of_linearEquiv}
        \leanok
  Write \(A = \Gamma(S, V)\) and \(R = \Gamma(S, U)\). Since \(M\) is flat over \(A\), the
  base change \(R \otimes_A M\) is flat over \(R\) (\cref{lem:mathlib_baseChange_flat}).
  Because \(A \to R\) is a ring epimorphism, the canonical map \(R \otimes_A M \to M\) is an
  isomorphism of \(R\)-modules (\cref{lem:mathlib_tensorProduct_lid_epi}); transporting
  flatness across this isomorphism (\cref{lem:mathlib_flat_of_linearEquiv}) shows \(M\) is
  flat over \(R\).
\end{proof}

\begin{lemma}\leanok
[Per-piece base descent via a base-change identity]
  \label{lem:gf_flat_of_isBaseChange_id}
  \lean{AlgebraicGeometry.gf_flat_of_isBaseChange_id}
  \uses{lem:mathlib_flat_isBaseChange}
  Let \(A \to R\) be a ring map and \(M\) an \(R\)-module that, through the scalar tower
  \(A \to R \to M\), is flat over \(A\). If the identity \(M \to M\) exhibits the
  \(R\)-module \(M\) as the base change of the flat \(A\)-module \(M\) along \(A \to R\)
  (equivalently, the canonical map \(R \otimes_A M \to M\) is an isomorphism), then \(M\)
  is flat over \(R\).
\end{lemma}

\begin{proof}
  \uses{lem:mathlib_flat_isBaseChange}
  \leanok
  Apply the base-change flatness principle \cref{lem:mathlib_flat_isBaseChange} to the
  hypothesised identity \(M \to M\): the target of a base-change identification of a flat
  module is flat, so \(M\) is flat over \(R\).
\end{proof}

\begin{lemma}\leanok
[B2 — source-span flatness descent for section modules]
  \label{lem:gf_section_span_flat_descent}
  \lean{AlgebraicGeometry.gf_section_span_flat_descent}
  \uses{lem:qcoh_section_localization_basicOpen, lem:mathlib_flat_of_isLocalized_span}
  Let \(p : X \to S\) be a morphism, \(\F\) a quasi-coherent \(\OO_X\)-module, \(U\) an open
  of \(S\), and \(W \le p^{-1}(U)\) an affine open of \(X\); give \(\Gamma(\F, W)\) its
  \(\Gamma(S, U)\)-module structure through the restriction
  \(\Gamma(S, U) \to \Gamma(X, W)\). Suppose a finite family \(t \subseteq \Gamma(X, W)\)
  spans the unit ideal, \(\mathrm{Ideal.span}\,t = \top\), and that for each \(g \in t\) the
  basic-open section module \(\Gamma(\F, D(g))\) — the away-localization of
  \(\Gamma(\F, W)\) at \(g\) — is flat over \(\Gamma(S, U)\). Then \(\Gamma(\F, W)\) is flat
  over \(\Gamma(S, U)\).
\end{lemma}

\begin{proof}
  \uses{lem:qcoh_section_localization_basicOpen, lem:mathlib_flat_of_isLocalized_span}
  \leanok
  By the section-localization datum \cref{lem:qcoh_section_localization_basicOpen}, for each
  \(g \in t\) the restriction \(\Gamma(\F, W) \to \Gamma(\F, D(g))\) exhibits
  \(\Gamma(\F, D(g))\) as the away-localization of \(\Gamma(\F, W)\) at \(g\). The family \(t\) spans the unit ideal of the source ring
  \(\Gamma(X, W)\), and each localized piece \(\Gamma(\F, D(g))\) is flat over the base
  \(\Gamma(S, U)\) by hypothesis. The source-span flatness criterion
  \cref{lem:mathlib_flat_of_isLocalized_span}, applied with base \(\Gamma(S, U)\), source
  \(\Gamma(X, W)\), and these principal localizations, yields that \(\Gamma(\F, W)\) is flat
  over \(\Gamma(S, U)\).
\end{proof}

\begin{lemma}

\leanok
[Common basic opens of two affines form a basis of the overlap]
  \label{lem:gf_common_basicOpen_basis}
  \lean{AlgebraicGeometry.gf_common_basicOpen_basis}
  \uses{lem:gf_crossChart_basicOpen_eq}
  Let \(X\) be a scheme with affine opens \(W, W_i\) and a point \(x \in W \cap W_i\). Then there
  exist sections \(g \in \Gamma(X, W)\) and \(\bar g \in \Gamma(X, W_i)\) with
  \(x \in D(g)\), \(D(g) \le W \sqcap W_i\), \(D(\bar g) \le W \sqcap W_i\), and a common basic
  open \(X.\mathrm{basicOpen}\,g = X.\mathrm{basicOpen}\,\bar g\). The opens simultaneously basic
  in \(W\) and in \(W_i\) thus form a basis of the overlap.
\end{lemma}

\begin{proof}
  \uses{lem:gf_crossChart_basicOpen_eq}
  \leanok
  Pick a basic open \(D(b)\) of \(W_i\) with \(x \in D(b) \subseteq W \cap W_i\)
  (\(\texttt{IsAffineOpen.exists\_basicOpen\_le}\) in chart \(W_i\)), then a basic open \(D(g)\) of
  \(W\) with \(x \in D(g) \subseteq D(b)\) (in chart \(W\)); hence \(D(g) \le W \sqcap W_i\).
  Realise \(D(g)\) back as a basic open of \(W_i\): the restriction \(g|_{D(b)}\) lives in
  \(\Gamma(X, D(b)) \cong (\Gamma(X, W_i))_b\) (\(\texttt{IsAffineOpen.isLocalization\_basicOpen}\)),
  so write \(g|_{D(b)} = \bar g'/b^{n}\) by \(\texttt{IsLocalization.surj}\) and set
  \(\bar g := b \cdot \bar g'\). Then \(X.\mathrm{basicOpen}\,\bar g = D(b) \sqcap X.\mathrm{basicOpen}\,\bar g' = X.\mathrm{basicOpen}\,g\)
  (\(\texttt{Scheme.basicOpen\_mul}\) plus \(\texttt{basicOpen\_res}\) bookkeeping, since \(b\) is a
  unit on \(D(b)\)), which is the basic-open equality. Note (iter-054 finding): only the basic-open
  equality is achievable in general; a restriction-matched pair \(g|_O = \bar g|_O\) need not exist.
\end{proof}

\begin{lemma}\leanok
[B2.4 — finite spanning basic-open cover aligned to the patches]
  \label{lem:gf_crossChart_spanning_cover}
  \lean{AlgebraicGeometry.gf_crossChart_spanning_cover}
  \uses{lem:gf_affine_finite_standard_subcover, lem:gf_common_basicOpen_basis, lem:gf_crossChart_basicOpen_eq}
  % NOTE (iter-054): the "matched pair … that agree over the overlap" (restriction-match
  % g|_O = ḡ|_O) is OVER-SPECIFIED — the Lean conclusion delivers ONLY the basic-open equality
  % `X.basicOpen g = X.basicOpen ḡ` (≤ W ⊓ Wⱼ). Prover finding: the restriction-matched pair is
  % NOT constructible in general (cross-chart realisation only yields ḡ = unit·g on the overlap);
  % the basic-open equality is what B2.2/the assembly consume. Planner: drop the "agree over the
  % overlap" clause here and in the proof sketch. Same correction applies to the (currently
  % unblueprinted) iter-054 helper `gf_common_basicOpen_basis`, which this lemma delegates to.
  Let \(X\) be a scheme, \(W \subseteq X\) an affine open, and \(\{W_j\}_{j}\) a finite family of
  affine opens of \(X\) whose union contains \(W\). Then there is a finite index set \(K\) and, for
  each \(k \in K\), a chart index \(j(k)\) together with sections \(g_k \in \Gamma(X, W)\),
  \(\bar g_k \in \Gamma(X, W_{j(k)})\) cutting out a common basic open
  \[
    D(g_k) \;=\; X.\mathrm{basicOpen}\,\bar g_k \;\le\; W \sqcap W_{j(k)},
  \]
  such that \(\mathrm{Ideal.span}\,\{g_k : k \in K\} = \top\) in \(\Gamma(X, W)\) (equivalently, the
  \(D(g_k)\) cover \(W\)). This is the spanning family fed to
  \cref{lem:mathlib_flat_of_isLocalized_span}.
\end{lemma}

\begin{proof}
  \uses{lem:gf_affine_finite_standard_subcover, lem:gf_crossChart_basicOpen_eq}
  \leanok
  Inside each overlap \(W \cap W_j\) — open in both affines \(W\) and \(W_j\) — the opens that are
  simultaneously basic in \(W\) and basic in \(W_j\) form a basis of the subspace topology: given
  \(x \in W \cap W_j\), choose \(b \in \Gamma(X, W_j)\) with \(x \in D(b) \subseteq W \cap W_j\) and
  then \(g \in \Gamma(X, W)\) with \(x \in D(g) \subseteq D(b)\); shrinking \(b\) to the common open
  and recording its restriction over the overlap produces a basic open of \(W\) that is matched, over
  \(W \cap W_j\), by a section \(\bar g \in \Gamma(X, W_j)\) (\cref{lem:gf_crossChart_basicOpen_eq}
  certifies \(D(g) = X.\mathrm{basicOpen}\,\bar g\)). Ranging over all \(j\), and because the \(W_j\)
  cover \(W\), these common basic opens form an open cover of the affine \(W\). The affine \(W\) is
  quasi-compact, so a finite subfamily already covers it: applying the standard-subcover extraction
  \cref{lem:gf_affine_finite_standard_subcover} to this open cover of \(W\) returns a finite family
  \(\{g_k\} \subseteq \Gamma(X, W)\) with \(\mathrm{Ideal.span}\,\{g_k\} = \top\), each \(D(g_k)\) a
  common basic open contained in some \(W \cap W_{j(k)}\) with matching
  \(\bar g_k \in \Gamma(X, W_{j(k)})\).
\end{proof}



\begin{theorem}



\leanok
[Generic flatness, geometric form]
  \label{thm:generic_flatness}
  \lean{AlgebraicGeometry.genericFlatness}
  \uses{thm:generic_flatness_algebraic, lem:isLocalizedModule_powers_restrictScalars}
  % SOURCE: [Nitsure], \S4, Theorem ``generic flatness''
  % (read from references/nitsure-hilbert-quot-src/nitsure-hilbert-quot.tex, L1781-L1787).
  % SOURCE QUOTE:
  % "Let \(S\) be a noetherian and integral scheme.
  %  Let \(p: X\to S\) be a finite type morphism, and let \(\F\) be a
  %  coherent sheaf of \({\mathcal O}_X\)-modules. Then there exists a non-empty
  %  open subscheme \(U\subset S\) such that the restriction of \(\F\)
  %  to \(X_U = p^{-1}(U)\) is flat over \({\mathcal O}_U\)."
  \textit{Source: [Nitsure], \S4, Theorem ``generic flatness''.}
  Let \(S\) be a noetherian integral scheme, \(p : X \to S\) a finite-type morphism,
  and \(\F\) a coherent \(\OO_X\)-module. Then there exists a non-empty open
  subscheme \(U \subseteq S\) such that the restriction of \(\F\) to
  \(X_U = p^{-1}(U)\) is \(\OO_U\)-flat.

  The coherence hypothesis on \(\F\) is essential: a non-coherent \(\OO_X\)-module
  need not become flat over any non-empty open of \(S\).
\end{theorem}

% SOURCE QUOTE PROOF: "The above theorem has the following consequence, which
% follows by restricting attention to a non-empty affine open subscheme of \(S\)."
% (Source: nitsure-hilbert-quot-src/nitsure-hilbert-quot.tex, L1776-L1777. The
% Nitsure proof is the single sentence quoted: deduce the geometric form from
% the algebraic generic flatness statement (\cref{thm:generic_flatness_algebraic})
% by restricting to a non-empty
% affine open of \(S\), where the situation reduces to the algebraic statement.)
\begin{proof}
  \uses{thm:generic_flatness_algebraic, lem:gf_qcoh_fintype_finite_sections,
        lem:gf_crossChart_spanning_cover, lem:gf_section_span_flat_descent,
        lem:gf_section_localization_twoleg, lem:gf_patch_free_imp_flat,
        lem:gf_flat_localizedModule_sameBase, lem:gf_openImmersion_isEpi,
        lem:gf_flat_descend_isEpi, lem:mathlib_flat_of_isLocalized_span,
        lem:mathlib_flat_of_free, lem:mathlib_free_of_isLocalizedModule,
        lem:gf_flat_isLocalizedModule_sameBase, lem:flat_of_ringEquiv_semilinear,
        lem:flat_localization_models, lem:isLocalizedModule_powers_restrictScalars}
        \leanok
  Restrict to a non-empty affine open \(U_0 = \Spec A \subseteq S\), where \(A\) is a
  noetherian domain (\(S\) integral and locally noetherian). We assemble the witness
  \(V\) from the algebraic form on a finite affine cover.

  \emph{Step 1 — finite affine cover above \(U_0\).} Since \(p\) is quasi-compact and
  \(U_0\) is affine, hence quasi-compact, the preimage \(X_{U_0} = p^{-1}(U_0)\) is
  quasi-compact; choose a
  finite affine cover \(\{W_j = \Spec B_j\}_{j}\) of \(X_{U_0}\) by affine opens of
  \(X\). Because \(p\) is locally of finite type, each ring \(B_j = \Gamma(X, W_j)\)
  is a finite-type \(A\)-algebra.

  \emph{Step 2 — read off a finite module on each patch.} On each \(W_j\), quasi-coherence
  of \(\F\) identifies \(\F|_{W_j}\) with the associated sheaf \(\widetilde{M_j}\) of the
  section module \(M_j = \Gamma(\F, W_j)\), and the finite-section bridge
  \cref{lem:gf_qcoh_fintype_finite_sections} — applied to the affine open \(W_j\) — makes
  \(M_j\) a finite \(B_j\)-module. Through the scalar tower \(A \to B_j \to M_j\),
  \(M_j\) is the data fed to the algebraic form.

  \emph{Step 3 — apply the algebraic form and take the common basic open.}
  \cref{thm:generic_flatness_algebraic} applied to \((A, B_j, M_j)\) yields \(f_j \in A\),
  \(f_j \ne 0\), with \((M_j)_{f_j}\) free over \(A_{f_j}\). As the cover is finite, set
  \(f := \prod_j f_j \ne 0\) (a non-zero element of the domain \(A\)) and
  \(V := D(f) \subseteq \Spec A \subseteq S\), a non-empty basic open; over \(V\) every
  \((M_j)_f\) is free over \(A_f\).

  \emph{Step 4 — freeness implies flatness (source-span descent with base descent along
  \(U \hookrightarrow V\)).} On each patch the localized section module \((M_j)_f\) is free,
  hence flat over \(A_f = \Gamma(S, V)\). We promote these per-patch freeness statements to
  the flatness conclusion for every affine \(U \le V\) and affine \(W \le p^{-1}(U)\), with
  \(R := \Gamma(S, U)\), by a \emph{source-span} descent on \(W\).

  Cover \(W\) by basic opens aligned to the patches: by the patch-aligned covering lemma
  \cref{lem:gf_crossChart_spanning_cover} there is a finite family
  \(\{g\} \subseteq \Gamma(X, W)\) with \(\mathrm{Ideal.span}\,\{g\} = \top\), each
  \(D(g) \subseteq W \cap W_{j}\) a common basic open matched by a patch element
  \(\bar g \in B_{j}\) (so \(D(g) = X.\mathrm{basicOpen}\,\bar g\)). For each such \(g\): the
  two-leg transport \cref{lem:gf_section_localization_twoleg} identifies the section module
  \(\Gamma(\F, D(g)) \cong (M_{j})_{\bar g}\). The per-patch freeness gives \((M_{j})_f\)
  free, hence flat over \(A_f\) (\cref{lem:gf_patch_free_imp_flat},
  \cref{lem:mathlib_flat_of_free}); since \(D(g)\) lies over \(U \le V = D(f)\) the image of
  \(f\) is already invertible on \(D(\bar g)\), so localizing \((M_{j})_f\) at the
  powers of \(\bar g\) keeps \(A_f\)-flatness
  (\cref{lem:gf_flat_localizedModule_sameBase}) and yields exactly \(\Gamma(\F, D(g))\) flat
  over \(A_f = \Gamma(S, V)\).

  \emph{Base descent along the open immersion.} The inclusion \(U \hookrightarrow V\) of
  affine opens makes \(\Gamma(S, V) \to \Gamma(S, U)\) a ring epimorphism
  (\cref{lem:gf_openImmersion_isEpi}); flat descent along a ring epimorphism
  \cref{lem:gf_flat_descend_isEpi} therefore upgrades ``\(\Gamma(\F, D(g))\) flat over
  \(\Gamma(S, V)\)'' to ``\(\Gamma(\F, D(g))\) flat over \(R = \Gamma(S, U)\)''. Finally the
  \(\{g\}\) span the unit ideal of \(\Gamma(X, W)\), so the source-span flatness criterion
  \cref{lem:mathlib_flat_of_isLocalized_span}, packaged as
  \cref{lem:gf_section_span_flat_descent}, descends these flat principal localizations to
  \(\Gamma(\F, W)\) flat over \(R = \Gamma(S, U)\). This holds for every affine \(U \le V\)
  and \(W \le p^{-1}(U)\), producing the required \(V\) with \(\F|_{X_V}\) flat over
  \(\OO_V\).
\end{proof}
